\begin{abstract}
We promote the $n$-birational motivic homotopy category $S\mapsto \mathcal{H}^n(S)$ to a $Pr^L$-valued presheaf on $Corr(\mathrm{Sch})_{uglt,sm}$. As a consequence, for any unibranch scheme $X$, its (zeroth) birational motivic homotopy category $\mathcal{H}^{b\mathbb{A}^1}(X)$ decomposes as the product of the birational motivic homotopy categories of its function fields; in particular, for a variety $V$, $\mathcal{H}^{b\mathbb{A}^1}(V) \simeq \mathcal{H}^{b\mathbb{A}^1}(k(V))$. This implies that birational equivalences of unibranch schemes in $Sm_S$ can be detected by the birational contractibility of their generic fibers. Finally, we establish that stably birational morphisms and purely transcendental field extensions induce fully faithful embeddings of birational motivic homotopy categories.
\end{abstract}
\tableofcontents

\section{Introduction}
\subsection{History of the subject and the aim of the paper}

There are several reasons to consider the birational localization of the motivic homotopy category. The initial and probably the most significant motivation for its development was the slice filtration (see [\cite{pelaez2014unstable}]). Another reason, evident from its identification with the birational homotopy category [\cite{0bat}, Corollary 2.2.12.(2)], is that it serves as the natural homotopical extension of classical birational geometry (see \S1.1. [\cite{0bat}]). (There is also a distinct reason: to incorporate topos-theoretic properties into the motivic homotopy category (see [\cite{inhot}]).)


Although most literature has studied the construction of the birational (motivic) homotopy category only over fields (see [\cite{choudhury2022characterisation}, [\cite{cisinski2023homotopy}]), one can easily extend its definition to general base schemes ([\cite{0bat}]). Namely, for a Qcqs scheme $S$, one defines the birational motivic homotopy category $\mathcal{H}^{b\mathbb{A}^1}(S)$ as the localization of the motivic homotopy category $\mathcal{H}^{\mathbb{A}^1}(S)$ at the set of all dense open immersions in $Sm_S$. 

Unfortunately, this unstable construction does not behave well under base change. The main issue is that the denseness of open immersions is not preserved when pulling back along arbitrary maps, particularly closed immersions (see the first paragraph of \S\ref{3.2} for an example). Additionally, this new localization is not locally cartesian, as shown in [\cite{0bat}, Counterexample 3.2.13]. As a result, the existence of an exceptional pullback along closed immersions is unclear.


Despite this drawback, we will prove in this paper that, for sufficiently well-behaved schemes, the associated birational motivic homotopy category can be assembled from its values at the generic points of the schemes (see \Cref{Hb of a var is Hb of its generic point}). 

To handle this elegantly, we introduce a more refined notion of $ n$-birational motives. Rather than localizing at all dense open immersions, we localize only at those whose closed complement has codimension greater than $n$. We call this the $n$-birational motivic category, denoted $\mathcal{H}^n(S)$. (Clearly, $\mathcal{H}^{b\mathbb{A}^1}(S)\equiv\mathcal{H}^0(S)$.)

This generalized notion first appeared in [\cite{pelaez2014unstable}] as the unstable counterpart to the stable slice category. The modern $\infty$-categorical avatars of these categories first appeared in the joint work of Bachmann and Elmanto [\cite{bachmann2019voevodsky}], where the authors used them to prove Voevodsky's slice conjectures [\cite{voevodsky2002open}]. Let us take a moment to explain why these are the correct unstable analogs of the stable notion of slices of motivic spectra. To simplify notation, we will fix a base Qcqs scheme $S$ for the subsequent discussion and functionally omit the symbol $S$ from the diagrams.


To set the stage, let us first recall the definition of motivic slices from [\cite{voevodsky2002possible}]. The full subcategory $\mathcal{SH}^{eff}(n)$ of $\mathcal{SH}_{\mathbb{P}^1}$, monoidally generated by $T^{\wedge n}$ and $\mathcal{SH}^{eff}$, is closed under arbitrary coproducts by construction. It is therefore, colocalizing due to compact generation and hence has exact colocalization functors, say given by $f_n:\mathcal{SH}_{\mathbb{P}^1}\to \mathcal{SH}^{eff}(n)$. The chain of inclusions $\mathcal{SH}^{eff}(n+1)\subset \mathcal{SH}^{eff}(n)$, upon colocalization, then induces a tower:
\begin{flalign}\label{slice tower}
    \cdots f_{n+1}\to f_n\to f_{n-1}\to \cdots 
\end{flalign}
called the slice tower [\S2, \cite{voevodsky2002open}]. The slices of this tower are called the motivic slices $s_n:=f_n/f_{n+1}$. 

One can perform a similar construction for the $S^1$-stable motivic homotopy category $\mathcal{SH}_{S^1}$. Namely, for every $n\geq 0$ consider the monoidal subcategory $\mathcal{SH}_{S^1}(n)$ of $\mathcal{SH}_{S^1}$ generated by $T^{\wedge n}$. In this case, one can similarly define a towoer of colization functors as in \Cref{slice tower} (still denoted by the same notation):
$$\cdots\to f_n\to f_{n-1}\to \cdots \to f_1\to f_0=1,$$ with slices denoted by $s_n=f_n/f_{n+1}$.

Due to the slice conjectures, the computation of positive slices in $\mathcal{SH}_{\mathbb{P}^1}$ can be computed under $\Omega^\infty_{\mathbb{G}}$ inside $\mathcal{SH}_{S^1}$ itself. It can be hoped that, in turn, the computation of slices in $\mathcal{SH}_{S^1}$ can be deduced from a similar construction in the unstable category. To do so effortlessly, one must, however, choose an appropriately designed unstable category. That is, for every $-_n=f_n, f_{0/n}, s_n$ one has to choose $\mathcal{H}^{-_n}$ such that $Stab(\mathcal{H}^{-_n})\simeq \mathcal{SH}^{-_n}_{S^1}$, where  $$\mathcal{SH}_{S^1}^{f_n}:=\mathcal{SH}_{S^1}(n),\text{ }\mathcal{SH}_{S^1}^{f_{0/n}}:=\mathcal{SH}_{S^1}/f_n\text{ or } \mathcal{SH}_{S^1}^{s_n}:=s_n\mathcal{SH}_{S^1}.$$ 

Let us now discuss the situation of a similar construction within the unstable setup. One quick observation is that, if one defines the subcategory $i_n:\mathcal{H}^{\mathbb{A}^1}_\bullet\wedge T^n\subset \mathcal{H}^{\mathbb{A}^1}_\bullet$ as the one monoidally generated by $T^{\wedge n}$, then it is closed under filtered colimits in $\mathcal{H}_\bullet^{\mathbb{A}^1}$. Indeed, if $$D:I\to \mathcal{H}^{\mathbb{A}^1}_\bullet \wedge T^n\xhookrightarrow{i_n}\mathcal{H}^{\mathbb{A}^1}_\bullet$$ is a filtered diagram, then, on the one hand, $\Sigma_T^n\Omega^n_Ti_nD\simeq D$ (by the triangle identity of the adjunction $\Sigma_T^n\dashv \Omega^n_T$), while, on the other hand, the endo-functor $\Sigma_T^n\Omega_T^n:\mathcal{H}^{\mathbb{A}^1}\to \mathcal{H}^{\mathbb{A}^1}$ preserves filtered colimits (indeed, $\Sigma_T^n$ is a left adjoint and thus preserves colimits, while the endo-functor $\Omega_T^n:\mathcal{H}^{\mathbb{A}^1}_\bullet\to \mathcal{H}^{\mathbb{A}^1}_\bullet$ preserves filtered colimits because $T$ ($\simeq \mathbb{P}^1$) is a compact object). Thus, $$\colim{i_nD}\simeq \colim\Sigma_T^n\Omega_T^n{i_nD}\simeq \Sigma_T^n\Omega_T^n\colim{i_nD}\in \mathcal{H}^{\mathbb{A}^1}_\bullet \wedge T^n.$$ That is, $i_n$ preserves filtered colimits. However, $i_n$ fails to preserve cofibers in general\footnote{For example, over a field $k$, consider the Hopf map $\eta : \mathbb{A}^2\setminus 0\to \mathbb{P}^1$. This has motivic cofiber identical to $\mathbb{P}^2$ [\cite{dugger2005motivic}, Proposition 2.13]. Even though $\mathbb{A}^2\setminus 0\simeq S^{3,2}\simeq T\wedge\mathbb{G}_m\in \mathcal{H}^{\mathbb{A}^1}_\bullet\wedge T\ni T\simeq \mathbb{P}^1$, we claim that $\mathbb{P}^2\notin \mathcal{H}^{\mathbb{A}^1}_\bullet\wedge T$. Suppose $\mathbb{P}^2\simeq^{mot} T\wedge E\simeq^{mot} \Sigma F$ for some spaces $E,F$. We find that $\mathcal{P}_{nis}(\Sigma F, \mathbb{G}_m)=*$, while $\mathcal{P}_{nis}(\mathbb{P}^2, \mathbb{G}_m)=k^\times$, and $\mathbb{G}_m$ is $\mathbb{A}^1$-local.}. This makes it impossible to ensure that it is colocalizing, unlike its stable counterpart. Hence, producing the correct unstable slice tower via right adjoints is not possible this way. What one usually does in such scenarios is to consider the colocalizing subcategory generated by the projection maps $T^n\to *$, rather than requiring $\mathcal{H}^{\mathbb{A}^1}_\bullet\wedge T^n$ to be colocalizing. The problem with this approach is that, in general, it is extremely hard to characterize colocalizing classes of morphisms. 

The problem is that, in the stable case, the aforementioned tower is the analog of the `Whitehead tower' of Tate connective covers (not the `Postnikov tower'). The Postnikov tower is actually the dual of this and can be identified (say in the $S^1$-case) with the layers of the above tower along $f_0$. In the unstable world, this is equivalent to studying the localizing subcategory induced by the projection $T^n\to *$, where the localizing class of morphisms is better handled. This approach has been used by many, most notably in [\cite{asok2023p}, Section 3], in greater generality of $S^{p,q}$ localization (though with a different goal of obtaining a Frudenthal $\mathbb{P}^1$-suspension theorem). It is evident that for any $p\geq q$ there is a canonical equivalence $Stab(L^{p,q}\mathcal{H}^{\mathbb{A}^1})\simeq \mathcal{SH}_{S^1}/f_q$. In fact, because of the equivalence $T^n\simeq \mathbb{A}^{n+1}\setminus 0$ [\cite{morel19991}, Example 3.2.20], one has $$Stab(L_{\mathbb{A}^{n+1}\setminus 0}\mathcal{H}^{\mathbb{A}^1})\simeq \mathcal{SH}_{S^1}/f_n.$$

No doubt, this model is the most canonical analog of the $S^1$ stable slice construction. However, deciding whether a morphism of schemes is an $S^{p,q}$ equivalence is a classic and quite hard problem. This is exactly where the birational construction comes into play: one makes the key observation that, over perfect fields, an $S^1$-spectrum is a $0$-slice (i.e., has trivial $1$-connective cover) if and only if it is local with respect to all dense open immersions [\cite{levine2008homotopy}, Corollary 4.1.5]. In fact, it follows that an $S^1$-spectrum has a trivial $(n+1)$-connective slice cover if and only if it is local with respect to all dense open immersions whose complements have codimension strictly greater than $n$ (call such morphisms $n$-dense, and $B_n$ their collection) [\cite{pelaez2013birational}]. That is, $\mathcal{SH}_{S^1}[B_n^{-1}]\simeq \mathcal{SH}_{S^1}/f_{n+1}$. 

One is then tempted to define the required unstable analogue of the stable slice categories $\mathcal{SH}_{S^1}/f_{n+1}$ as the localization of $\mathcal{H}^{\mathbb{A}^1}$ at the set of $n$-dense open immersions, denoted by $\mathcal{H}^{n}$, so as to obtain an equivalence $$Stab(\mathcal{H}^{n})\simeq \mathcal{SH}_{S^1}/f_{n+1},$$ but this time over perfect fields. These $\mathcal{H}^{n}$'s are exactly the unstable $n$-birational motivic homotopy categories we mentioned in the third paragraph. Note that, because of classical birational geometry, it is easier to detect morphisms in the corresponding birational saturation classes than in the $(\mathbb{A}^{n+1}\setminus0)$-localized saturation class.

While the slice categories appear to have good functoriality properties under base change, it is understandable that the birational construction will not behave as well as the original motivic homotopy construction. To elaborate, recall that the six-functor formalism for the stable motivic homotopy category is first developed in pieces at the unstable level and then glued at the stable level using the localization sequence. Let us now compare the unstable birational construction to the original motivic homotopy construction from this functorial viewpoint.


As a first example, recall that the motivic homotopy construction $\mathcal{H}^{\mathbb{A}^1}$ is readily a presheaf on $\mathrm{Sch}$ taking values in $Pr^L$. However, as we have mentioned earlier, because $n$-dense maps are not stable under arbitrary pullbacks, the schematic functoriality of $\mathcal{H}^{n}$ (when taking values in $Pr^L$) on all morphisms of $\mathrm{Sch}$ is not immediate. Morphisms in $\mathrm{Sch}$ along which $\mathcal{H}^n$ has this functorial behavior will be those that can lift dense open immersions. In this paper, we will identify one such class, denoted $UGLT$ (see \Cref{uglt defn}).

Another key functorial property of the unstable motivic construction is that the presheaf $\mathcal{H}^{\mathbb{A}^1}: \mathrm{Sch}\to Cat_\infty$ satisfies Nisnevich descent [\cite{hoyois2017six}, Proposition 4.8]. In this paper, we will prove that the same holds for the unstable $n$-birational motivic homotopy construction when restricted to an appropriate domain that ensures functoriality.

It is likewise predictable that some of our results will be weaker than their $\mathbb{A}^1$ motivic counterparts. A notable example is the idea of essentially smooth base change, which, in the ordinary motivic case, says that $\mathcal{H}^{\mathbb{A}^1}$ is continuous with respect to pro-systems with affine transition maps. This turns out to be a very powerful tool in motivic homotopy for reasons worth mentioning. First, the Noetherian decomposition of qcqs schemes suggests that it suffices to study $\mathcal{H}^{\mathbb{A}^1}$ over Noetherian schemes alone, as was primarily done by the pioneers themselves [\cite{morel19991}]. Another important consequence is that, when working over fields, it (mostly) suffices to assume that the field is perfect (since every field is a filtered colimit of affine schemes over a perfect field). Again, for such a continuity result to work birationally, the pro-system must preserve dense immersions under base change. Unfortunately, neither of the examples of pro-systems mentioned earlier is of this type. Nonetheless, we will prove an appropriate statement for $\mathcal{H}^n$ and see that, even with this limited continuity, it yields meaningful results for the $0$-th category $\mathcal{H}^0\equiv\mathcal{H}^{b\mathbb{A}^1}= \mathcal{H}^b$ [\cite{0bat}].

The intent of this paper can thus be summarized as follows: 
\begin{enumerate}
    \item to determine the correct functoriality of $\mathcal{H}^n$;
    \item to prove the appropriate relations among those functors;
\end{enumerate}

Let us outline some implications of these for the category of our primary interest, the birational motivic homotopy category $\mathcal{H}^0\equiv\mathcal{H}^{b\mathbb{A}^1}$ [\cite{0bat}]. We will show that $\mathcal{H}^{b\mathbb{A}^1}$ is a birational invariant (induces an equivalence along birational maps), $\mathbb{A}^1$-invariant (induces a fully faithful embedding along $\mathbb{A}^1$-projections), and deformation invariant (induces an equivalence along nil immersions) Nisnevich sheaf of $\infty$-categories satisfying flat continuity. In particular, we demonstrate that the birational homotopy category of a suitably nice (unibranch) scheme is a direct product of the birational homotopy categories of its generic points. Using these insights, we characterize birational motivic equivalences of such well-behaved smooth schemes over a base.

\subsection{Notations and terminologies}
Throughout this paper, all schemes are Qcqs. By a morphism of schemes, we shall always mean a separated morphism. Smooth morphisms will be assumed to be quasicompact and hence of finite presentation. By a birational morphism, we mean a span $X\hookleftarrow
 U \hookrightarrow Y$ of dense open immersions. $X/S$ is said to be rational if it is birational to $\mathbb{A}^n_S$ (over $S$). $\mathrm{Sch}$ denotes the category of all qcqs schemes. $\mathrm{Noet}$ and $\mathrm{Nuc}$ denote the full subcategories of $\mathrm{Sch}$ consisting of Noetherian, respectively Noetherian and universally catenary schemes. 

We shall freely use the language of $\infty$-categories without referring to any particular model for the theory. The reader may freely use their preferred model. A standard reference is [\cite{lurie2009higher}], where the author models these on Kan complexes as models for $\infty$-groupoids. With a choice of Grothendieck universes, the terms small, large, and very large will be treated as usual. $Pr^L$ shall denote the (very large) $\infty$-category of large presentable $\infty$-categories, while $Cat_{\infty}$ will denote the (very large) $\infty$-category of all large $\infty$-categories. $Pr^{L,\otimes}$ will denote the standard monoidal structure on $Pr^{L,\otimes}$ [\cite{lurie2017higher}, \S6.3]. We will use the notation $\mathrm{CAlg}(Pr^L)$ for the full subcategory of commutative algebra objects of $Pr^L$.


For an $\infty$-category $\mathcal{C}$, we denote by $\mathcal{P}(\mathcal{C})$ the category of presheaves of spaces on $\mathcal{C}$ and by $h$ the Yoneda embedding $\mathcal{C}\to \mathcal{P}(\mathcal{C})$. If $\mathcal{C}$ has coproducts, we denote by $\mathcal{P}_\Sigma(\mathcal{C})$ the full subcategory of $\mathcal{P}(\mathcal{C})$ consisting of presheaves that turn finite coproducts into products. For a composition-closed class $M\subset Mor(\mathcal{C})$, we denote by $\mathcal{C}^M$ the wide subcategory spanned by morphisms in $M$. For an object $S\in \mathcal{C}$, we denote by $\mathrm{M}_S$ the full subcategory of $\mathcal{C}_{/S}$ spanned by $M$-morphisms over $S$. Finally, we denote by $\mathrm{M}^M_S$ the category $\mathcal{C}^M_{/S}$.



Given a Qcqs scheme $S$, we denote by $Sm_S$ the full subcategory of $\mathrm{Sch}_{/S}$ consisting of schemes over $S$ that are relatively smooth and finitely presented. While one might consider the larger category $SM_S$ of (not necessarily finitely presented) smooth schemes over $S$, it turns out that Nisnevich locally this is irrelevant, in that one has an equivalence $\mathcal{P}_{nis}(Sm_S)\simeq \mathcal{P}_{nis}(Sm_S^{fp})$ [\cite{hoyois2015quadratic}, Appendix C, Proposition C.5 (3)]. Because we always work Nisnevich locally, we shall work with the (essentially) small category $Sm_S$. By $\mathcal{P}(S)$, we shall mean the $\infty$-category of spaces on $Sm_S$, i.e., $\mathcal{P}(S):=Fun(Sm_S^{op},Spc)$. For a smooth scheme $X/S$, we denote by $h_SX$ the corresponding representable presheaf.

We will use the notation $\mathcal{H}^{\mathbb{A}^1}(S)$ for the full subcategory of $\mathcal{P}(S)$ consisting of $\mathbb{A}^1$-local (i.e., local with respect to the set of projection maps $\{X\times \mathbb{A}^1\to X\}_{X\in Sm_S}$) Nisnevich sheaves of spaces. Equivalently, $\mathcal{H}^{\mathbb{A}^1}(S)=L_{\mathbb{A}^1}L_{nis}\mathcal{P}(S)$. The associated localization functor is denoted by $L_{mot}$. Morphisms in the associated saturated class will be called Motivic equivalences. 
\pagebreak

\subsection{Details and motivations}
Given a (Qcqs) scheme $S$, we define the $n$-Birational Motivic homotopy category $\mathcal{H}^{n}(S)$ as the localization of the motivic homotopy category $\mathcal{H}^{\mathbb{A}^1}(S)$ at the set of dense open subschemes whose complements have codimension strictly greater than $n$ (we call such morphisms $n$-dense open immersions). Since $Sm_S$ is essentially small, this localization is accessible [\cite{lurie2009higher}] and is given by an accessible localization functor, which we denote by $L^n$. 

Let us recall that, although the $\mathbb{A}^1$ motivic localization functor is not entirely left exact, it has several good properties that make ordinary motivic homotopy theory at least manageable. For example, $L_{mot}$ is both cartesian and locally cartesian. These two properties hold because the $\mathbb{A}^1$ projection maps are stable under arbitrary pullbacks (and hence under products with smooth schemes). The situation with the $n$-Birational localization is a little disappointing, in that, since dense open immersions are not stable under arbitrary pullbacks, it is not clear whether the $n$-Birational localization functor is locally cartesian. However, it is cartesian for smooth morphisms. In fact, these problems regarding the failure of pullback stability of dense open immersions under arbitrary pullbacks cause another serious drawback: the pullback-pushforward adjunction does not descend to the level of the $n$-birational homotopy category. In fact, these new localizations are quite strange from a functorial perspective.

To elaborate on this, let us try to extend the definition of the $n$-Birational motivic homotopy category construction to a `presheaf' $\mathcal{H}^{n}(-)^*$ of $\infty$-categories on $\mathrm{Sch}$, by assigning $f\mapsto L_{n}f^*b^n$, where $b^n$ denotes the inclusion $\mathcal{H}^n\subset \mathcal{P}_{nis}$. The problem with this construction is that, for it to preserve compositions, $f^*$ must preserve $n$-Birational equivalences. Unfortunately, this does not hold for all morphisms. A good class of morphisms that can do this job is the collection of those that universally lift generalizations [\Cref{ugltpull}]. This yields the presheaf $\mathcal{H}^{n}(-)^*:Sch^{uglt,op}\to Pr^L$. Obviously, the smooth extension functor $p_\shrp$, induced by a smooth map $p$, trivially preserves $n$-Birational equivalences, giving rise to the co-presheaf $\mathcal{H}^n(-)_\shrp: Sch^{sm}\to Pr^L$ mapping $p\mapsto L^np_\shrp$. This is the broad content of \S3.1 and can be summarized as saying:
\begin{recall*}{full functoriality}
   There is a presheaf $$\mathcal{H}^n(-): Corr(\mathrm{Sch})_{sm,uglt}^{op}\to Pr^L$$ whose restriction to the vertical maps recovers $\mathcal{H}^n(-)_\shrp: \mathrm{Sch}^{sm}\to Pr^L$, while the restriction to the horizontal maps recovers $\mathcal{H}^{n}(-)^*: \mathrm{Sch}^{uglt,op}_{}\to Pr^L$.
\end{recall*}

The basic problem with full functoriality of $\mathcal{H}^n$ is largely tied to non-flat closed immersions. Specifically, when pulling back a dense open to a closed subset, the intersection may become empty, say by missing some generic points of the closed subscheme, if not all of them. This matter is delicate and requires careful treatment. Instead of asking for stability of arbitrary dense open immersions under closed pullback, we ask for a dimension-dependent version of the same, and show in \Cref{closed rel d push} that for Noetherian Universally Catenary schemes, the graded assignment $\bigsqcup \mathcal{H}^n(-): Ob(\mathrm{Nuc})\to Gr_\mathbb{N}Cat_\infty$ extends to a \textit{codimension}-graded presheaf $$\mathcal{H}^\mathbb{N}(-)^*:\mathrm{Nuc}^{equi-closed,op}\to Pr^L\text{ given by }i\mapsto L^{n}i^*: \mathcal{H}^{{n+mcod(i)}}(S) \rightarrow \mathcal{H}^{n}(Z)$$ 


The first surprising yet naively expected result in this section is the following:

    \begin{recall*}{dense equiv}
       For a dense open immersion $j:V\xhookrightarrow{}S$, the smooth extension map $$j_{0\shrp}:=L_{0}j_\shrp :\mathcal{H}^0(V)\to \mathcal{H}^0(S)$$ is an equivalence, and hence so are $j_*\simeq L_{0}j_*$ and $j^*\simeq L_{0}j^*$.
    \end{recall*}
We use the term `surprise' to emphasize that this is not a general phenomenon. To elaborate, suppose we study the $P$-invariant theory inside a suitable topos $\mathcal{P}_\tau(Sm_{\mbox{-}})$ (conceptually, this is obtained by localizing at a class of morphisms in $Sm_{\mbox{-}}$ with property $P$). In general, there is no reason to believe that the presheaf $\mathcal{H}_{\shrp}^{\{P\}}: \mathrm{Sch}^{sm}\to Pr^L$ is itself $P$-invariant. The $\mathbb{A}^1$ homotopy category is an immediate example (see the paragraph above \Cref{motivic category is a1 faithful}). 

On this line, one might expect that the philosophy of the main observation in [\cite{0bat}], that "`whatever' is birational local is $\mathbb{A}^1$ local", should apply immediately to this scenario and yield that the birational homotopy presheaf is itself $\mathbb{A}^1$-invariant, in the strongest sense of inducing a categorical equivalence against $\mathbb{A}^1$ projections. That is, the dense local invariant presheaf $\mathcal{H}^{0}(-)^*\equiv \mathcal{H}^{b\mathbb{A}^1}(-)^*$, with values in $Pr^L$ as above, is also $\mathbb{A}^1$-invariant (unlike the $\mathbb{A}^1$ homotopy category, which induces only a fully faithful embedding under the $\mathbb{A}^1$ projections). The first problem is that if we work with the presheaf $$\mathcal{H}^0(-)^*\equiv \mathcal{H}^{b\mathbb{A}^1}(-)^*: \mathrm{Sch}^{uglt,op}\to Pr^L,$$ we miss an important ingredient of the original argument in [\cite{cisinski2023homotopy}]. Namely, their proof of the $\mathbb{A}^1$-invariance of a birational presheaf involves certain blow-ups that are not necessarily glt (or smooth). At the end of \S4.3, we discuss in detail why the arguments in [\cite{cisinski2023homotopy}, Theorem 1.7] are not applicable to the birational presheaf $\mathcal{H}^0(-)^*: \mathrm{Sch}^{uglt,op}\to Pr^L$. However, we do have:

\begin{inner-recall-star}[Homotopy invariance, see \Cref{homotopy}]
    For a scheme $S$, the pullback maps $p^{n*}: \mathcal{H}^{n}(S)\to \mathcal{H}^{n}(\mathbb{A}^m_S)$ along the projection $p:\mathbb{A}^m_S\to S$ are fully faithful.
\end{inner-recall-star}



Let us also recall another surprising result about birational presheaves. It states that additive presheaves that are dense local are automatically Nisnevich local [\cite{0bat}, Theorem 2.1.2]. It is therefore natural to expect that our $\mathcal{H}^{0}(-)^*$ will satisfy Nisnevich descent. For this, however, it must belong to $\mathcal{P}_\Sigma (-,Pr^L)$ (see [\cite{0bat}, Theorem 2.1.2]; or \Cref{nis is bir} in this paper). That is, it must transform disjoint unions of schemes into products of $\infty$-categories. While it is easy to prove the last claim and thus settle the case of Nisnevich descent for $\mathcal{H}^0$, it turns out that the descent property can be proved by hand simultaneously for all $n$ and all schemes. The point is that the Nisnevich topos is itself Nisnevich local [\cite{hoyois2017six}, Proposition 4.8], so it remains to observe that a collection of compatible dense open immersions given over a Nisnevich cover descends `uniquely' to the total space of the cover:
 \begin{inner-recall-star}[Nisnevich descent, see \Cref{bir detect}]
  The presheaf $$
\mathcal{H}^{n}(-)^*: \mathrm{Sch}^{uglt,op}\to Pr^L$$ is a Nisnevich sheaf.       
    \end{inner-recall-star}
There are a few more continuity properties enjoyed by our $n$-Birational motivic homotopy construction. For example, we show that,

\begin{inner-recall-star}[{Deformation-invariant}, see \Cref{deformation invariant}]
The presheaf $$\mathcal{H}^{n}(-)^*: Sch^{uglt,op}\to Pr^L$$ sends thickenings of schemes to equivalences of $\infty$-categories. Specifically, if $Z\xhookrightarrow{i} S$ is a nil immersion, then $i_*:  \mathcal{H}^{n}(Z) \to \mathcal{H}^{n}(S)$ is an equivalence of $\infty$-categories. In particular, the canonical map $$\mathcal{H}^{n}(X^{red})\to \mathcal{H}^{n}(X)$$ given by pullback along the reduction $X^{red}\subset X$ is an equivalence.
    \end{inner-recall-star}
As expected, we obtain a flat affine continuity result for $\mathcal{H}^0$ (see \Cref{continuity of Hb}). Combined with the above results, we then conclude that,
  \begin{inner-recall-star}[Generic decomposition, see \Cref{Hb of a var is Hb of its generic point}]
         Let $X$ be a Qcqs unibranch scheme. Then the canonical map $$\mathcal{H}^0(X)\to\underset{\eta\in X^{(0)}}{\prod}\mathcal{H}^0(k(\eta))$$ is an equivalence. In particular, for a variety $V$, $\mathcal{H}^0(V)\simeq\mathcal{H}^0(K(V))$.
     \end{inner-recall-star}
As an immediate corollary, we obtain a fiberwise criterion for smooth $0$-birational motivic equivalences of schemes:
\begin{inner-recall-star}[Fibersise criteria for birational equivalences, see \Cref{fibersise criteria}]
    Suppose $S$ is a Qcqs geometrically unibranch scheme and $f:X\to Y$ is a morphism in $Sm_S$ such that the generic fibers of $f$ are birationally contractible, i.e., for every generic point $\eta\in Y^{(0)}$ the scheme $f_\eta:X_{k(\eta)}\to k(\eta)$ is contractible in $\mathcal{H}^0(k(\eta))$. Then $f$ is a birational equivalence in $\mathcal{P}(S)$.
\end{inner-recall-star} 
Additionally, the generic decomposition (\Cref{Hb of a var is Hb of its generic point}) yields:
\begin{inner-recall-star}[Rational invariance, see \Cref{rational invariance}]
Let $S$ be a scheme and assume that $X$ is a rational scheme over $S$ (in other words, $X\to S$ is a stably birational isomorphism). Then there is a canonical embedding $\mathcal{H}^0(S)\hookrightarrow \mathcal{H}^0(X)$.
    \end{inner-recall-star}
    
\begin{inner-recall-star}[Pure transcendence invariance, see \Cref{trans}]
   Let $k$ be a field. Then the embedding $k\subset k(t_1,\cdots,t_n) $ induces an embedding $\mathcal{H}^0(k)\subset \mathcal{H}^0(k(t_1,\cdots,t_n))$.
\end{inner-recall-star}
 
We end this paper in \S\ref{4.5} with a proof of a partial version of the unstable slice conjecture:

\begin{recall*}{coneavue adjunction}
     For every $n,d\geq 0$
      $$\Omega^n _{\mathbb{p}^1}\mathcal{H}^{{n+d}}_\bullet(S) \subset\mathcal{H}^{{d}}_\bullet(S).$$
  \end{recall*}
This reduces the unstable version of the slice conjecture (i.e., $\Omega_{\mathbb{P}^1}L^{n+1}\simeq L^n\Omega_{\mathbb{P}^1}$) to showing that:
\begin{conj}[see \Cref{conj}]
For a Qcqs scheme $S$, the canonical map $L^n\Omega_{\mathbb{P}^1}\to \Omega_{\mathbb{P}^1}L^{n+1}$ induced by \Cref{coneavue adjunction}, is an $n$-birational equivalence (on simply connected motivic spaces; see \Cref{Spq slice conj} for the justification of the condition).
\end{conj}



\textit{\textbf{What we do not do in this paper: }}
\begin{enumerate}
    \item {We will only set up the unstable theory here; the stable theory, as one always does in a paper on the six-functor formalism, could follow with a little amount of functorial work. In our case, however, this might require more correction to achieve a full six-functor formalism, and I postpone it to my thesis.} 
    \item The other issue is that the $n$-birational motivic homotopy construction is not locally cartesian (see [\cite{0bat}, Counterexample 3.2.13] for $n=0$). Consequently, many motivic functoriality results that use this property fail to descend to the $n$-birational level. For example, we do not have an exceptional pushforward functor for closed immersions.
\end{enumerate}
\section{Recollections on \texorpdfstring{$n$}{}-Birational Motivic homotopy theory}
\subsection{\texorpdfstring{$n$}{}-Birational geometry}
Let $X$ be a scheme. If $x,y\in X$ are points, then $y$ is a generalization of $x$, written $x\leftsquigarrow y$, if $x\neq y$ and $x\in \bar{y}$. A chain of generalizations starting at $x$ (or, equivalently, a generic chain of $x$) is a sequence $$x_0\leftsquigarrow x_1\leftsquigarrow x_2\leftsquigarrow \cdots \leftsquigarrow x_n$$ of consecutive generalizations with $x_0=x$. For such a chain, $n$ is called its length.
 
When $x\in X$, we denote by $cod_Xx$ the supremum of the lengths of generalizing chains starting at $x$. In other words,  $cod_Xx:=codim(\bar{x},X)$ (see [\cite[\href{https://stacks.math.columbia.edu/tag/02I3}{Definition 02I3}]{stacks-project}] for notation). From [\cite[\href{https://stacks.math.columbia.edu/tag/02IZ}{Lemma 02IZ}]{stacks-project}] it follows that $cod_Xx=dim(\mathcal{O}_{X,x})$, where the right-hand side denotes the Krull dimension of the local ring $\mathcal{O}_{X,x}$. Given a natural number $d$, $X^{(d)}$ will stand for the set of all codimension $d$ points of $X$.


More generally, for a closed subscheme $i:Z\subset X$, we denote $$cod{(i)}: =\underset{\eta_Z\in Z^{(0)}}{min} \{cod_X(\eta_Z)\}.$$ We will also denote it by $cod_X(Z)$ when the immersion $i$ is clear from the context. We shall also use the notation $$mcod{(i)}: =\underset{\eta_Z\in Z^{(0)}}{sup} \{cod_X(\eta_Z)\}$$ and $\delta (i)=mcod(i)-cod(i)$. We say that $i$ is equidimensional if $\delta(i)=0$. 


Now we come to the most important geometric notion of this paper:

\begin{defn}
An open immersion $j: U\hookrightarrow X$ of schemes is said to be $n$-dense if $cod_X(X\setminus j U) >n$. 
\end{defn}

In other words, an open immersion $U\hookrightarrow X$ is $n$-dense if every generic point (and hence every other point) of $X\setminus jU$ has codimension at least $n+1$ in $X$, or, equivalently, if every codimension $\leq n$ point of $X$ belongs to $U$. In [\cite{bachmann2019voevodsky}, Definition 2.2], this notion is called an `$n$-birational morphism'. We call these $n$-dense to emphasize that we are not yet dealing with a span.

Since a generalization need not lift along all morphisms, $n$-dense open immersions are not stable under arbitrary pullbacks. In fact, lifting generalizations seems to be sufficient for this purpose:
\begin{lem}
    Suppose $f: S\to T$ lifts generalizations [\cite[\href{https://stacks.math.columbia.edu/tag/0063}{Definition 0063}]{stacks-project}]. Let $U\subset T$ be an $n$-dense open immersion. Then $f^{-1}U\subset S$ is also $n$-dense.
\end{lem}

\begin{proof}
To see this, let $s\in S-f^{-1}(U)$. This implies that $f(s)\in T-U$, and thus has codimension at least $n+1$. Any maximal chain of generalizations starting at $f(s)$ must therefore have length at least $n+1$. Since $f$ lifts generalizations, we may lift this chain to obtain a chain of generalizations starting at $s$ and thus of length at least $n+1$. This implies that $cod_S(S-f^{-1}(U))\geq n+1$. 
\end{proof}

\subsection{\texorpdfstring{$n$}{}-Birational Motivic spaces}
In this section, we will recall the construction of $n$-Birational motivic homotopy theory from [\cite{bachmann2019voevodsky}]. To set the stage, we will first recall the construction of the ordinary motivic homotopy category $\mathcal{H}^{\mathbb{A}^1}(S)$ for an arbitrary (Qcqs) scheme $S$. 

We denote by $Sm_S$ the full subcategory of $\mathrm{Sch}_S$ consisting of schemes over $S$ with quasi-compact, separated, and smooth structure maps. Since $S$ is Qcqs and the objects of $Sm_S$ are qcqs morphisms, $Sm_S$ is an (essentially) small (1-)category. We denote by $\mathcal{P}(S)$ the $\infty$-category of presheaves of spaces on $Sm_S$, i.e., $\mathcal{P}(S):=Fun(Sm_S^{op}, Spc)$. We call these $S$-spaces.

A morphism of schemes $X\to Y$ in $Sm_S$ is a Nisnevich cover if it is an \'etale cover (i.e., a surjective \'etale map) and is completely decomposed (i.e., induces at least one isomorphism of residue fields) at each point of $Y$. A space $\esc{X}\in \mathcal{P}(S)$ is a Nisnevich sheaf (or Nisnevich local) if it satisfies descent with respect to every Nisnevich cover. We denote the full subcategory of $\mathcal{P}(S)$ consisting of Nisnevich sheaves by $\mathcal{P}_{nis}(S)$. Since $Sm_S$ is (essentially) small, it follows that $\mathcal{P}_{nis}(S)\subset \mathcal{P}(S)$ is a left-exact, accessible localization, which we denote by $L_{nis}$, called the Nisnevich sheafification functor. A morphism of $\mathcal{P}(S)$ inverted by $L_{nis}$ is called a Nisnevich local equivalence, and we denote the collection of such morphisms by $nis(S)$. We denote this new $\infty$-category by $\mathcal{H}^{\mathbb{A}^1}(S)$.

Next, we call an object $\esc{X}\in \mathcal{P}(S)$ to be $\mathbb{A}^1$-local if it is local with respect to the class $\mathbb{A}(S):=\{\mathbb{A}^1_S\times X\to X\}$ of $\mathbb{A}^1$-projection maps. The motivic homotopy category is the full subcategory of the $\infty$-topos of Nisnevich sheaves consisting of $\mathbb{A}^1$-local objects. Its objects, i.e., those presheaves of spaces on $Sm_S$ that are both Nisnevich-local and $\mathbb{A}^1$-local, are called motivic spaces. 

The definition of the $n$-Birational Motivic homotopy categories is similar. For each $n\in \mathbb{N}$, we let $$B_n(S):=\{U\xhookrightarrow{i} X \in Mor(Sm_S)\mid cod_X(X-iU)\geq n+1\}$$ (that is, $B_n(S)$ is the set of $n$-dense open immersions contained in $Sm_S$). We say that a Motivic space $\esc{X}\in\mathcal{H}^{\mathbb{A}^1}(S)$ is an $n$-birational motivic space if it is, moreover, local with respect to the set $B_n(S)$. In other words:
\begin{defn}
    An $S$-space $\esc{X}$ is called an $n$-birational motivic space if it is a Nisnevich sheaf and is local with respect to the class $\mathbb{A}(S)\cup B_n(S)$ of $\mathbb{A}^1$-projections and $n$-dense open immersions in $Sm_S$. We denote by $\mathcal{H} ^{n}(S)$ the full subcategory of $\mathcal{P}(S)$ consisting of $n$-birational motivic spaces.
\end{defn}

By definition, we have a natural chain of full subcategories $$\mathcal{H}^n(S)\subset \mathcal{H}^{\mathbb{A}^1}(S)\subset \mathcal{P}_{nis}(S)\subset \mathcal{P}(S).$$ Since $\mathbb{A}(S)\cup B_n(S)\subset Mor(Sm_S)$ and $Sm_S$ is (essentially) small, it follows that both $\mathbb{A}(S)$ and $B_n(S)$ are (essentially) small. Thus, by [\cite{lurie2009higher}, Proposition 5.5.4.15], each possible composition in the chain of inclusions $$\mathcal{H}^n(S)\subset \mathcal{H}^{\mathbb{A}^1}(S)\subset \mathcal{P}_{nis}(S)\subset \mathcal{P}(S)$$ is an accessible localization. We denote the localization functor for $\mathcal{H}^{\mathbb{A}^1}(S)\subset \mathcal{P}(S)$ by $L_{mot}$ and call a morphism in $\mathcal{P}(S)$ inverted by $L_{mot}$ a motivic equivalence, and denote the collection of such morphisms by $mot(S)$. Similarly, we denote the localization functors for both the inclusions $\mathcal{H}^n(S)\subset \mathcal{H}^{\mathbb{A}^1}(S)$ and $\mathcal{H}^n(S)\subset \mathcal{P}(S)$ by $L^n_S$ and call a morphism in $\mathcal{P}(S)$ inverted by $L^n_S$ an $n$-Birational motivic equivalence, and denote the collection of such morphisms by $bir_n(S)$. 
\begin{rem}
    Since an open immersion is isomorphic to the open subscheme of its image, the \textit{set} of open subschemes forms a skeleton of $B_0(S)$ and generates the same localizing subcategory. We therefore treat "dense open immersion" and "dense open subscheme" interchangeably in the subsequent results.
\end{rem}

It follows from [\cite{lurie2009higher}, Proposition 5.5.4.15] that, in each case, the classes of morphisms inverted by the localization functors are precisely the saturation classes generated by (i.e., the smallest class closed under pushouts, two-out-of-three, and small colimits [\cite{lurie2009higher}, Definition 5.5.4.5] and containing) the class of morphisms being localized. For example, $mot(S)$ is the saturation class generated by Nisnevich equivalences and the set $\mathbb{A}(S)$, while $bir_n(S)$ is generated by motivic equivalences and $B_n(S)$, or equivalently by Nisnevich equivalences and $\mathbb{A}(S)\cup B_n(S)$.

In fact, both subcategories, $L_{\mathbb{A}^1}\mathcal{P}(S)$ and $L_{dense}^n\mathcal{P}(S)$, of the $\infty$-topos $\mathcal{P}(S)$ (consisting of $\mathbb{A}^1$-local and $B_n(S)$-local objects, respectively) are accessible localizations. We denote the corresponding localization functors by $L_{\mathbb{A}^1}$ and $L_{dense}^n$. Elements of the naive saturation classes generated by $\mathbb{A}(S)$ and $B_n(S)$ will be called homotopy equivalences and $n$-dense equivalences, respectively. It follows that $mot(S)$ is the saturation class generated by $nis(S)$ and homotopy equivalences; similarly, $bir_n(S)$ is the saturation class generated by $mot(S)$ and $n$-dense equivalences. We will drop $S$ from all notations when either the base scheme is clear from the context or we are working functionally over all schemes. 


By the same proposition [\cite{lurie2009higher}, Proposition 5.5.4.15], both $\mathcal{H}^{n}(S)$ and $ \mathcal{H}^{\mathbb{A^1}}(S)$ are presentable $\infty$-categories. The candidate $\mathcal{H}^{\mathbb{A^1}}(S)$ first appeared in [\cite{morel19991}] in model-categorical terms and has since been studied extensively by various authors. The other candidate $\mathcal{H}^n(S)$ first appeared in modern $\infty$-categorical language in [\cite{bachmann2019voevodsky}], where it was defined for locally noetherian schemes and denoted $L^n_{bir}\mathcal{H}^{\mathbb{A}^1}(S)$, with the localization functors denoted $L^n_{bir}$. Since, in this paper we will always be working with birational localizations, we have removed $bir$ from the subscript to simplify notation. When $n=0$, the candidate $\mathcal{H}^0(S)$ is identical to $\mathcal{H}^{b\mathbb{A}^1}(S)$ of [\cite{0bat}], where it is shown to coincide with the birational homotopy category $\mathcal{H}^b(S):=L_{dense}\mathcal{P}_\Sigma(S)$ (see Corollary 2.2.10 (2) of \textit{loc.cit.}) (we will reprove this claim here using somewhat different arguments, see \Cref{Motivic equivalences are Birational equivalences}). Finally (to complete the analogy), we recognize $\mathcal{H}^{{-1}}(S)\simeq \{\underline{\emptyset}\to*\}=:\mathcal{I}$, the interval category, and $L^{-1}\simeq  \tau_{\leq -1}^{nis}$.

Because the contractible $S$-space is trivially local, it is also a terminal object of $\mathcal{H}^n(S)$. The pointed category $\mathcal{H}_\bullet ^n(S)$ is the slice category $\mathcal{H}^n(S)_{*/}$. By [\cite{lurie2009higher}, 5.5.2.10], this is also a presentable $\infty$-category. More specifically, the forgetful functor $\mathcal{H}_\bullet^n(S)\to \mathcal{H}^n(S)$ has a left adjoint, characterized by preservation of colimits and by $h_S(X)\mapsto h_S(X)_+$, where this last term is the space $h_S(X\sqcup S)$ given by the obvious inclusion $S\to X\sqcup S$. We denote this left adjoint, which adds base points, by $(-)_\bullet$. By the general machinery of $Pr^{L,\otimes }$, we have $\mathcal{H}^n_\bullet(S)\simeq \mathcal{H}^n(S)\otimes Spc_*$, where $Spc_*$ is the category of pointed spaces.

\textbf{The birational slice tower}\label{bst}

Clearly, $B_{n+1}(S)\subset B_n(S)$, so $\mathcal{H}^{n}(S)\subset \mathcal{H}^{{n+1}}(S)$. This immediately yields a chain of inclusions of $\infty$-categories:
$$\mathcal{I}\simeq \mathcal{H}^{{-1}}(S)\subset \mathcal{H}^{b}(S)=\mathcal{H}^{0}(S)\subset \mathcal{H}^{1}(S)\subset \mathcal{H}^{2}(S)\subset \cdots \subset \mathcal{H}^{\infty}(S)\subset\mathcal{H}^{\mathbb{A}^1}(S) $$
(here, $\mathcal{H}^\infty(S)$ is the colimit of all the previous inclusions) or, equivalently, a chain of inclusions among the saturation classes:
$$all(S)\supset bir_0(S)\supset bir_1(S)\supset bir_2(S)\supset \cdots \supset bir_\infty(S)\supset mot(S) $$
That is, all $n$-Birational motivic equivalences are, in particular, $0$-birational motivic equivalences. We choose to call a $0$-birational motivic equivalence simply a Birational motivic equivalence. In this terminology, it makes sense to call an $n$-Birational motivic equivalence a Birational motivic equivalence of height $n$.


By denoting the relative localizations $\mathcal{H}^{n+1}(S)\to \mathcal{H}^{n}(S)$ (corresponding to the inclusion $\mathcal{H}^n(S)\subset \mathcal{H}^{n+1}(S)$) still by $L^n$, the above chains yield a tower in $Pr^L$:
$$\mathcal{I}\simeq \mathcal{H}^{{-1}}(S)\xleftarrow{L_{}^{-1}}\mathcal{H}^{0}(S)\xleftarrow{L_{}^{0}} \mathcal{H}^{1}(S)\xleftarrow{L_{}^{1}}\mathcal{H}^{2}(S)\xleftarrow{L_{}^{2}} \cdots \leftarrow\mathcal{H}^{\infty}(S)\xleftarrow{L^{\infty}_{}}\mathcal{H}^{\mathbb{A}^1}(S) $$
We call this the unstable Birational slice tower. The pointed categories will be denoted by $\mathcal{H}^n_\bullet (S)$ and can similarly be assembled into a tower in $Pr^L$:
$$*\simeq \mathcal{H}_\bullet^{{-1}}(S)\xleftarrow{L_{}^{-1}}\mathcal{H}_\bullet^{0}(S)\xleftarrow{L_{}^{0}} \mathcal{H}_\bullet^{1}(S)\xleftarrow{L_{}^{1}}\mathcal{H}_\bullet^{2}(S)\xleftarrow{L_{}^{2}} \cdots \leftarrow\mathcal{H}_\bullet^{\infty}(S)\xleftarrow{L^{\infty}_{}}\mathcal{H}_\bullet^{\mathbb{A}^1}(S) $$
\begin{rem}\label[rem]{2.2.3}
    It is unclear whether these towers converge to the ordinary motivic homotopy category, i.e., whether the induced morphism in $Pr^L$ $$\mathcal{H}^{\infty}(S):=\plim_{n}\mathcal{H}^n(S)\xleftarrow{L^{\infty}}\mathcal{H}^{\mathbb{A}^1}(S) $$ is an equivalence of $\infty$-categories. In particular, suppose we are given a tower of spaces $$\esc{X}_0\xleftarrow{p_0}\esc{X}_1\xleftarrow{p_1}\esc{X}_2\xleftarrow{p_2}\cdots$$ with each $\esc{X}_i\in \mathcal{H}^{i}(S)$ and such that each $p_i$ is an $i$-birational motivic equivalence (equivalently, the induced morphisms $L^{i+1}\esc{X}_{i+1}\to \esc{X}_{i}$ are equivalences) for all $i\geq0$. Then it is not clear whether there is a motivic space $\esc{X}\in \mathcal{H}^{\mathbb{A}^1}(S)$ with a compatible family of $i$-birational motivic equivalences $\pi_i:\esc{X}\to \esc{X}_i$ (i.e., such that the induced morphisms $L^i\esc{X}\to \esc{X}_i$ are equivalences).
\end{rem}
\subsection{On equivalent constructions of $\mathcal{H}^n$}
In this section, we discuss some equivalent ways the $n$-birational homotopy category could have been constructed.

\underline{{Changing the base category from $Sm_S$}}

It turns out that the Nisnevich topos $\mathcal{P}_{nis}(S)$ can also be constructed from other subcategories of schemes over $S$. Let us recall the general mechanism for doing so. Let $AffSm_S$ denote the full subcategory of $Sm_S$ consisting of smooth $X/S$ such that $X=SpecR$ is an affine scheme. Let $a: AffSm_S\hookrightarrow Sm_S$ be the obvious inclusion. Since it is continuous and co-continuous for the Nisnevich topology, there is an adjunction of the sheaf $\infty$-toposes with a fully faithful right adjoint:
$$a^*:\mathcal{P}_{nis}(Sm_S)\leftrightarrows \mathcal{P}_{nis}(AffSm_S):a_*$$

When $S$ is separated, every Smooth (separated) $X/S$ can be covered by objects of $(AffSm_S)_{/X}$ whose finite intersections lie in $AffSm_S$. From standard \v{C}ech nerve computations, this immediately induces an equivalence of the sheaf toposes ([\cite{MR3679884} Lemma 3.3.2]; see also [\cite[\href{https://stacks.math.columbia.edu/tag/03A0}{Lemma 03A0}]{stacks-project}], though written for discrete spaces, applies equally to spaces):
$$a^*:\mathcal{P}_{nis}(Sm_S)\simeq\mathcal{P}_{nis}(AffSm_S):a_*$$
More generally, when $S$ is semiseparated, the functor $a$ factors through the full subcategory $SmAff_S$ of $Sm_S$ consisting of smooth affine morphisms $X/S$, which in turn factors through the subcategory $SmQproj_S$ of smooth quasiprojective morphisms $X/S$. It is evident that all these inclusions are continuous and cocontinuous for the Nisnevich topology. Therefore, when $S$ is separated (so that we can always choose a Zariski cover of any smooth $X/S$ such that all finite intersections are affine), the above equivalence factors as a composition of equivalences:
$$\mathcal{P}_{nis}(Sm_S)\simeq\mathcal{P}_{nis}(SmQproj_S)\simeq  \mathcal{P}_{nis}(SmAff_S)\simeq  \mathcal{P}_{nis}(AffSm_S)$$

In fact, when $S$ is separated, so that, so is every $X$ for $X/S\in Sm_S$, in the Nisnevich topos $\mathcal{P}_{nis}(Sm_S)$, we may write $X$ as a colimit of objects of $AffSm_S$, and thus every $\mathbb{A}^1$ projection map in $Sm_S$ (thought of in $\mathcal{P}_{nis}(S)$) is a colimit of $\mathbb{A}^1$ projections in $AffSm_S$ [\cite{MR3679884}, Lemma 5.1.2]. It follows that the above equivalences descend to equivalences of the ordinary motivic homotopy categories:
\begin{prop}
    Suppose $S$ is a separated scheme. Then there are equivalences of presentable $\infty$-categories:
$$\mathcal{H}^{\mathbb{A}^1}(Sm_S)\simeq\mathcal{H}^{\mathbb{A}^1}(SmQproj_S)\simeq  \mathcal{H}^{\mathbb{A}^1}(SmAff_S)\simeq  \mathcal{H}^{\mathbb{A}^1}(AffSm_S)$$
\end{prop}
We now discuss the $n$-birational case. The same methods apply immediately to this case. Indeed, if $U_i$ is a Zariski cover of $X$, then the open immersions $U_i\subset X$ lift generalizations, so that for any $n$-dense open $V\hookrightarrow X$, the base change $V\times _XU_i\hookrightarrow U_i$ is $n$-dense, and $V\times _XU_i$ forms a Zariski cover of $V$. Thus, we obtain the following result, where the notation $\mathcal{H}^n(\mathcal{C}_S)$ stands for $L_{dense}^nL_{\mathbb{A}^1}L_{nis}\mathcal{P}(\mathcal{C}_S)$ for an (admissible) subcategory $\mathcal{C}_S\subset Sm_S$:
\begin{thm}\label{smqproj}
 Let $S$ be a separated scheme. Then, for every $n$, there is a canonical equivalence of presentable $\infty$-categories:
   $$\mathcal{H}^{n}(Sm_S)\simeq\mathcal{H}^{n}(SmQproj_S)\simeq  \mathcal{H}^{n}(SmAff_S)\simeq  \mathcal{H}^{n}(AffSm_S)$$ 
\end{thm}

\textbf{Reducing local conditions for constructing $\mathcal{H}^0$}\label{Hb=H0}

It turns out that when $n=0$, the corresponding birational motivic homotopy category can be constructed in many other ways. Some of these constructions are discussed in detail in another work of mine [\cite{0bat}]. For example, $\mathcal{H}^0$ is equivalent to the $0$-dense localization of the logarithmic motivic homotopy category of [\cite{binda2023logarithmic}], as shown in [\cite{0bat}, Corollary 2.3.6]. In this document, we will present different lines of argument (aligned with the content of this paper) for the main model of $\mathcal{H}^0$ from \textit{loc.cit.}. In that paper, the arguments were mostly from the perspective of local objects; here, instead, we shall work with local equivalences. 

Let us now recall the construction of the birational homotopy category, $\mathcal{H}^b$, from [\cite{0bat}]. For a scheme $S$, let $\mathcal{P}_\Sigma(S)$ be the full subcategory of $\mathcal{P}(S)$ consisting of presheaves that send finite disjoint unions of schemes to products of spaces. This is the sifted co-completion of $Sm_S$ [\cite{lurie2009higher}, Proposition 5.5.8.10]. Since coproducts in $Sm_S$ are disjoint, this is also the category of sheaves for the disjoint-union topology [\cite{bachmann2017norms}, Lemma 2.4]. Local equivalences for this topology are precisely the \v{C}ech nerves of morphisms of the following types $$\bigsqcup _{i=1}^kh_S{U_i}\to h_S({\bigsqcup _{i=1}^kU_i}).$$ 

Let $\mathcal{H}^b(S)$ denote the full subcategory of $\mathcal{P}_\Sigma(S)$ consisting of $0$-dense local objects; i.e., $\mathcal{H}^b(S):=L^0_{dense}\mathcal{P}_\Sigma(Sm_S)$. Denote the associated saturation class by $bir_\Sigma(S)$ and the corresponding localization functor by $L_{bir}$. We first show that $nis(S)\subset bir_\Sigma(S)$. For this, we shall need some preliminary lemmas:

\begin{lem}\label[lem]{nis dense splitting}
    Let $p: U\to X$ be a Nisnevich cover of qcqs schemes; then $p$ splits over a dense open subscheme $V\subset X$.
\end{lem}
\begin{proof}
    Since $X_{nis}=X_{nis}^{red}$ and $X^{red}\subset X$ is a nil immersion, we may assume that $X$ is reduced. Indeed, suppose that on the dense open $V^{red}\subset X^{red}$, the Nisnevich cover $U^{red}\to X^{red}$ admits a section. Since both $U^{red}$ and $V^{red}$ are étale over $X^{red}$, the canonical base change equivalence of categories $Et/X\simeq Et/{X^{red}}$ [\cite{grothendieck1966elements}, Théorème (18.1.2)] implies that there is an open $V\subset X$ with reduction $V^{red}$ and a section $V^{}\to U^{}$ over $X^{}$ reducing to the section $V^{red}\to U^{red}$ over $X^{red}$. Since $X^{red}\subset X$ is a homeomorphism, the denseness of $V\subset X$ follows from the denseness of the pullback $V^{red}\subset X^{red}$.
    
    Now, by [\cite{hoyois2016trivial}], there is a stratification of finitely presented closed subschemes $$\emptyset=Z_0\subset Z_1 \subset Z_2\dots\subset Z_{n-1}\subset Z_n=X$$ such that $p$ admits a section $p_i$ over each locally closed stratum $S_i=Z_i\setminus Z_{i-1}$. Let $\mathscr{I}_i$ be the finitely presented ideal defining $Z_i$. If $\eta$ is a generic point of $X$ contained in $S_{i_\eta}$, then reducedness implies $(\mathscr{I}_{i_\eta})_{\eta}=0$. Since $\mathscr{I}_{i_\eta}$ is coherent, there exists an open $V_{\eta}\subset X$ such that ${\mathscr{I}_{i_\eta}}_{|V_{\eta}}=0$. Define, $$W_\eta:=V_\eta\bigcap S_{\eta_i}=V_\eta\bigcap (X\setminus Z_{\eta_i-1}).$$ Clearly, $W_\eta$ is open in $X$, contains $\eta$, and is contained in $S_{i_\eta}$. 

    If $X^{(0)}$ is the set of generic points of $X$, then for each $i$ consider the open set $$Y_i:=\bigcup_{\eta\in{X^{(0)}\bigcap S_i}}W_\eta\subset S_i.$$ The restriction ${p_i}_{|Y_i}$ provides a section of $p$ over $Y_i$. Since the $S_i$'s are disjoint from one another, so are the $Y_i$'s. Hence, the open set $Y=\bigsqcup Y_i$ admits a section. By construction, $X^{(0)}\subset Y$. Hence, $Y$ is dense in $X$, and we are done. \end{proof}

\begin{lem}\label[lem]{dense prod}
    Let $X_i\to S$ be a finite family of maps, each universally lifting generalizations. Assume $U_i\xhookrightarrow{}X_i\in bir_{n_i}(S)$. Then $$\prod_S U_i\xhookrightarrow{} \prod_S X_i\in bir_ {min \{n_i\}}{(S)}.$$
\end{lem}
\begin{proof}
    First, observe that $bir_{n_i}(S)\subset {min \{n_i(S)\}}$, so we may work with a single ${n}$. By induction, we may further reduce to the case of a two-fold product. But even then $f\times g=(f\times1)\circ (1\times g)$. Thus, it suffices to prove that if $U\xhookrightarrow{} X$ is in $B_n(S)$ and $Y\to S$ has universal generalization lifting, then $U\times_S Y\to X\times_S Y$ is in $B_n(S)$ as well. But this is the pullback of $U\xhookrightarrow{} X$ along $X\times_S Y\to X$. Since the latter lifts generalization (due to universality), we are done by the \Cref{glt pull}.
\end{proof}


\begin{thm}[see \cite{0bat}, Theorem 2.1.3. for a different argument]\label{nis is bir}
    Let $S$ be a qcqs scheme. Every Nisnevich equivalence in $\mathcal{P}(Sm_S)$ lies in $bir_\Sigma(S)$.
\end{thm}
\begin{proof}
Since $S$ is qcqs, it suffices to show that the \v{C}ech nerve of a basic Nisnevich cover belongs to $bir_\Sigma(S)$ [see \cite{hoyois2017six}, Proposition 3.8]. 

Let $X/S$ be smooth and let $\{U_i\to X\}_{i=1}^{k}$ be a Nisnevich cover of $X$. Since (the \v Cech nerve of) $$\bigsqcup h_{U_i}\to h_{\bigsqcup _iU_i}$$ is an equivalence in $ \mathcal{P}_\Sigma(Sm_S)$, by taking a disjoint union $U=\bigsqcup_i U_i$ we are left to show that if $p:U\to X$ is a Nisnevich cover given by a single morphism of schemes, then the colimit of the \v{C}ech nerve of the morphism $h_p:h_SU\to h_SX$ is a $0$-dense equivalence.
    
  By \Cref{nis dense splitting}, there is a dense open subset $V\subset X$ and a section $V\to U$ of $p$. Consider the pullback square on the left and the commutative diagram on the right, consisting of the \v{C}ech nerves of the vertical arrows in the left square:
  \[\xymatrix{
  (U\times_XV) \ar@{^(->}[r]^-{j'}\ar[d]_{q} &U\ar[d]^{p}&&  \check{C}(q) \ar[r]^-{\check{C}(j')}\ar[d]_{q_\bullet} &\check{C} (p)\ar[d]^{p_\bullet}&&\colim\check{C}(q) \ar[r]^-{|\check{C}(j')|}\ar[d]_{|q_\bullet|} &\colim \check{C} (p)\ar[d]^{|p_\bullet|}\\
 V\ar@{^(->}[r]_{j}& X&&  h_SV\ar[r]_{h_j}& h_SX&&  h_SV\ar[r]_{h_Sj}& h_SX
  }\]
  
  As $p$ lifts generalizations, \Cref{dense prod} implies that the morphism of \v{C}ech nerves $\check{C}(q) \xrightarrow{\check{C}(j')}\check{C} (p)$ is dense in each degree and hence a simplicial $0$-dense equivalence. Since the class of $0$-dense equivalences (being saturated) is closed under colimits, it follows that the top horizontal map $|\check{C}(j')| $ in the right-hand square is a dense equivalence. On the other hand, since $q:U\times_XV\to V$ splits, the left vertical map in the right-hand square is an equivalence in $\mathcal{P}(Sm_S)$. Because the bottom map $h_Sj$ is a dense equivalence by construction, the 2-out-of-3 property of the class of dense equivalences, together with the commutativity of the right-hand square, yields the desired result that $|p_\bullet|$ is a dense equivalence.
\end{proof}  

\begin{rem}
    The proof above offers a refinement of the argument presented in [\cite{choudhury2022characterisation}, Lemma 3.6]. While their approach effectively addresses the case of a single Nisnevich cover, it requires modification to fully account for Nisnevich covers arising from disjoint union decompositions. Consequently, the specific notion of birational equivalences utilized there (which correspond to $0$-dense equivalences in the present work) does not encompass all Nisnevich equivalences. Indeed, as observed in the first paragraph of the proof above, for a nontrivial decomposition $A\bigsqcup B=X$ of smooth schemes over $k$, the \v{C}ech nerve of the presheaf map $h_SA\bigsqcup h_SB\to h_SX$ is not automatically a sectionwise equivalence in $\mathcal{P}(k)$.
\end{rem}
The following result shows that $\mathcal{H}^b$ is an equivalent model to $\mathcal{H}^0$:
\begin{cor}[\cite{0bat}, Corollary 2.2.9]\label{Motivic equivalences are Birational equivalences}
Suppose $S$ is a qcqs scheme. Then the inclusion $\mathcal{H}^0(S)\subset \mathcal{H}^b(S)$ induces an equivalence of $\infty$-categories.
\end{cor}
By \Cref{nis is bir}, the above theorem follows from the following lemma, whose proof appears in [\cite{MR3404383}, Appendix A] (the proof has also been recalled in [\cite{cisinski2023homotopy}, Proposition 1.3], which we prefer to use below for its diagrammatic presentation). We recall it here for further use and the convenience of the reader:
\begin{lem}\label[lem]{bir local is P1 local}
    Let $S$ be a (QCQS) scheme and $\mathcal{G}$ a Grothendieck universe. Then every $0$-dense local presheaf $F\in \mathcal{G}\mbox{-}PSh_S$ taking values in the category of $\mathcal{G}$-sets is $\mathbb{P}^1$-local, and hence also $\mathbb{A}^1$-local. 
\end{lem} 
\begin{proof}
It suffices to prove that the morphism $F(S)\to F(\mathbb{P}^1_S)$ induced by the projection $\mathbb{P}^1_S\to S$ is an equivalence. We suppress the base $S$ from the notation. As in [\cite{cisinski2023homotopy}, Proposition 1.3], consider the blowup $W\to \mathbb{P}^2$ of $\mathbb{P}^2$ at a pair of closed points, say $\{0,\infty\}$, and let $j:E\rightarrowtail W$ be the exceptional fiber over $0$. It turns out that $W$ is also the blowup of $\mathbb{P}^1\times \mathbb{P}^1$ at the closed point $(0,0)$, which normalizes the birational isomorphism $\mathbb{P}^1\times\mathbb{P}^1\dashrightarrow \mathbb{P}^2$. Moreover, $E$ maps isomorphically to one copy of $\mathbb{P}^1$. This situation is summarized by the following diagram on the left, while the right-hand diagram is obtained by applying the contravariant functor $F$ to the left-hand diagram:
\[
\xymatrix@C=1.75em{
  & E \ar[dl]_{\pi} \ar@{>->}[d]_{j} \ar[dr]^{\sim} & & & & & F(E) & & \\
  S \ar@{>->}[d]_i & W \ar[dl]_{a} \ar[dr]^{b} & L \ar@{>->}[d]_k \ar[dr]^{\sim} & \ar@{|->}[r]^F & & F(S) \ar[ur]^{F(\pi)} & F(W) \ar@{->>}[u]_{F(j)} & F(L) \ar[ul]_{\sim} & \\
  \mathbb{P}^2 & & \mathbb{P}^1 \times \mathbb{P}^1 \ar[r]^{p_1} & \mathbb{P}^1 & & F(\mathbb{P}^2) \ar[u]^{F(i)} \ar@{->>}[ur]_{F(a)} & & F(\mathbb{P}^1 \times \mathbb{P}^1) \ar[ul]^{F(b)} \ar@{->>}[u]^{F(k)} & F(\mathbb{P}^1) \ar[l]^-{F(p_1)} \ar[ul]_{\sim}
}
\]
Any contravariant functor $F$ takes $k$ to a surjection, since $Fk\circ Fp_1$ is an isomorphism. But $Fj\circ Fb$ is isomorphic to $Fk$, so $Fj$ must be surjective. Since $a$ is a span of $0$-dense open immersions, $Fa$ is a surjection (in fact, an isomorphism). Because $F\pi\circ Fi= Fj\circ Fa$ and both $Fj$ and $Fa$ are surjective, we deduce that $F\pi$ is surjective. Since $\pi$ is split, it follows that $F\pi$ is an isomorphism. But $E\cong \mathbb{P}^1$. Thus $FS\underset{F\pi}{\twoheadrightarrow}FE\cong F\mathbb{P}^1$ given by the projection $\mathbb{P}^1\to S$ is an isomorphism.
\end{proof}
\begin{rem}
    Since $a$ is not $n$-dense for any $n>0$, the method above does not apply to the general $n$-birational case. On the other hand, the proof applies to any presheaf that turns the blow-up $W\to \mathbb{P}^2$ into a surjection, not just $0$-dense local ones.
\end{rem}


\subsection{On generating $\mathcal{H}^n$}
Except for the $n=0$ cases, the following result is standard:
\begin{prop}\label[prop]{compact generation}
    \begin{enumerate}
        \item $\mathcal{H}^{n}(S)$ is generated under sifted colimits by $n$-birational models of smooth schemes $X/S$. When $S$ is separated, we may take $X$ to be affine (in the absolute sense). When $S$ is the spectrum of a field of characteristic zero, we may take $X$ to be smooth and proper.
        \item  Each inclusion in the chain $$\mathcal{H}^n(S)\subset \mathcal{P}_{nis}(Sm_S)\subset\mathcal{P}(Sm_S)$$ is closed under filtered colimits. When $n=0$ and $S$ is qcqs, each inclusion in the chain $$\mathcal{H}^0(S)\subset \mathcal{P}_\Sigma(Sm_S)\subset\mathcal{P}(Sm_S)$$ is, in fact, closed under sifted colimits. If $S$ is furthermore geometrically unibranch, then $\mathcal{H}^0(S)\subset \mathcal{P}_\Sigma(Sm_S)$ is closed under all colimits.
        \item Every smooth $X/S$ is a compact object in $\mathcal{H}^n(S)$. When $n=0$, these generators are projective (i.e., they preserve geometric realizations) as well.
    \end{enumerate}
\end{prop}     

\begin{proof}
   (1). By construction, $\mathcal{P}_{\Sigma}(Sm_S)$ is generated under sifted colimits by smooth $X/S$. Being a localizing subcategory, so is $\mathcal{H}^{n}(S)$. By choosing Zariski covers (note that Zariski hypercovers are birational), we may take the (absolute) affine schemes smooth over $S$ as the generators (see the discussions around \Cref{smqproj}). For the last statement of (1) see [\cite{0bat}, Lemma 3.1.7 (c)].

(2). Since the inclusion $\mathcal{P}_{nis}(Sm_S)\subset\mathcal{P}(Sm_S)$ is closed under filtered colimits, while the inclusions $L_{\mathbb{A}^1}\mathcal{P}(Sm_S)\subset \mathcal{P}(Sm_S)$ and $L_{dense}\mathcal{P}(Sm_S)\subset \mathcal{P}(Sm_S)$ are closed under all colimits, the result follows immediately from the definition: $$\mathcal{H}^n(S)=L_{\mathbb{A}^1}\mathcal{P}(S)\bigcap L^n_{dense}\mathcal{P}(S)\bigcap \mathcal{P}_{\Sigma}(Sm_S).$$ When $n=0$ and $S$ is qcqs, we use the fact that $$\mathcal{H}^0(S)=\mathcal{H}^b(S)=L_{dense}\mathcal{P}(Sm_S)\bigcap \mathcal{P}_{\Sigma}(Sm_S)$$ (see \Cref{Motivic equivalences are Birational equivalences}) along with the observation that $\mathcal{P}_\Sigma(Sm_S)\subset\mathcal{P}(Sm_S)$ is closed under sifted colimits. The last statement follows from the additional observation that when $S$ is geometrically unbranch, the inclusion $\mathcal{H}^0(S)\subset \mathcal{P}_\Sigma(Sm_S)$ is closed under small coproducts [\cite{0bat}, Proposition 3.1.5].

(3). Every smooth $X/S$ is a compact projective object in $\mathcal{P}(Sm_S)$. Since projective objects are defined by preservation of geometric realizations and compact objects by preservation of filtered colimits, we are done by (2).
\end{proof}
Given a smooth scheme $X/S$ we shall use the notation  $h^n_S(X):=L^n_S(h_S(X))$

\begin{rem}
We caution the reader that none of the localizations $L^n$ is locally cartesian. The problem is again that $n$-dense open immersions are not stable under arbitrary pullback.
\end{rem}
\begin{cor}
Let $S$ be geometrically unibranch. Then $\mathcal{H}^0(S)\subset \mathcal{P}_\Sigma(Sm_S)$ is colocalizing.
\end{cor}
\begin{proof}
    Since both categories are presentable, using \Cref{compact generation}(2), this follows from the adjoint functor theorem [\cite{lurie2009higher}, Corollary 5.5.2.9(1)].
\end{proof}
\subsection{Monoidal structure on \texorpdfstring{$\mathcal{H}^n$}{hb} and \texorpdfstring{$\mathcal{H}^n_\bullet$}{hb} and their schematic functorialities}
We continue to work over a fixed Qcqs base scheme $S$. In this section, we will lift the pointed and unpointed birational motivic $\infty$-categories and their associated unstable functors to (symmetric) monoidal ones. Since the inclusion functor $\mathcal{H}^{n}(S)\subset \mathcal{H}^{\mathbb{A}^1}(S)$ preserves limits, it lifts immediately to a monoidal functor between the cartesian monoidal $\infty$-categories ${\mathcal{H}^{n}}(S)^\times\subset {\mathcal{H}^{\mathbb{A}^1}}(S)^\times$. However, this construction yields only an op-lax monoidal structure on the left adjoint $L^n:{\mathcal{H}^{\mathbb{A}^1}}(S)^\times\to {\mathcal{H}^{n}}(S)^\times$. That this is monoidal follows from the following theorem of Lurie and \Cref{bir cart}:
\begin{prop}[Proposition 1.3.9.;\cite{lurie2007derived}]\label{lurie 1.3.9}

Let $\mathcal{C}^{{{\otimes}}}$ be a monoidal structure on the $\infty$-category $\mathcal{C}$, and let $L : \mathcal{C} \to \mathcal{C}$ be a compatible localization functor with essential image $\mathcal{D} \subset \mathcal{C}$. Let $\mathcal{D}^\otimes$ be the full subcategory of $\mathcal{C}^\otimes$ spanned by objects in $\mathcal{D}$ (see the first paragraph of [\cite{lurie2007derived}, \S1.3]). Then:
\begin{enumerate}
\item $\mathcal{D}^{{\otimes}}$ equips the $\infty$-category $\mathcal{D}$ with a monoidal structure.
    \item The inclusion $\mathcal{D}^{{\otimes}} \subset \mathcal{C}^{{\otimes}}$ admits a left adjoint $L^{{\otimes}}$.
    \item The inclusion functor $\mathcal{D}^{{\otimes}} \subset \mathcal{C}^{{\otimes}}$ is lax monoidal, and its left adjoint $L^{{\otimes}} : \mathcal{C}^{{\otimes}} \to \mathcal{D}^{{\otimes}}$ is monoidal.
\end{enumerate} 
\end{prop} 
Recall that, following Remark 1.3.6 of \textit{loc.cit.}, for a closed (symmetric) monoidal presentable $\infty$-category $\mathcal{C}$, if the localization $L\mathcal{C}=\mathcal{D}$ is obtained by inverting $S\subset Mor\mathcal{C}$, then the condition that $L$ is compatible with the monoidal structure is equivalent to the statement that for any $f\in S$ and $Z\in \mathcal{C}$, $f\otimes Z\in S$. In our situation, we thus have the following:

\begin{prop}\label[prop]{bir cart}
    The birational localization $L^n$ is compatible with the cartesian monoidal structure of $\mathcal{P}_{nis}(Sm_S)$.
\end{prop}
\begin{proof}
    For a scheme $S$, recall that $\mathcal{P}_{nis}(S)$ is an $\infty$-topos, compactly generated by smooth schemes $X/S$. Thus, by the strong saturation of $bir_n(S)$, it suffices to verify that for any smooth scheme $p:X \to S$ and an open immersion $Y\xhookrightarrow{j} Z\in B_ n(S)$ of smooth schemes $Y/S$ and $Z/S$, the map $X\times_S Y\xhookrightarrow{X\times j}X\times_S Z$ lies in $B_n(S)$. But this is clear from \cref{dense prod}.
\end{proof}
Thus applying \Cref{lurie 1.3.9} we have:
\begin{thm}\label{Ln is cart}
    For every scheme $S$, the localization functor $L^n: \mathcal{P}(S)\to \mathcal{H}^{n}(S)$ lifts to a monoidal functor $\mathcal{P}(S)^{\times}\xrightarrow[]{L_{bir}^{n\times}} \mathcal{H}^{n}(S)^\times$. In fact, the entire chain of localization functors $$\mathcal{P}(S)\xrightarrow[]{L_{mot}} \mathcal{H}^{\mathbb{A}^1}(S)\xrightarrow[]{L^n} \mathcal{H}^{n}(S)$$ lifts to a corresponding chain of monoidal $\infty$-categories: $$\mathcal{P}(S)^{\times}\xrightarrow[]{L_{mot}^\times } \mathcal{H}^{\mathbb{A}^1}(S)^\times\xrightarrow[]{L^{n\times}} \mathcal{H}^{n}(S)^\times.$$
    \end{thm}
\begin{proof}
    For the first statement, combine \Cref{lurie 1.3.9} and \Cref{bir cart}. The rest is obvious because $L_{mot}$ is Cartesian [\cite{hoyois2017six}, Proposition 3.15.]. 
\end{proof}

Note that, by the above theorem, we are also saying that the induced monoidal structure on $\mathcal{H}^n(S)$ is cartesian. In particular, since the original cartesian structure is closed, we obtain an internal hom-tensor adjunction at the birational level: $$\esc{X}\times-:\mathcal{H}^{n}(S)\rightleftarrows\mathcal{H}^{n}(S): (-)^{\esc{X}}$$

\begin{rem}
    The categorical moral of the above machinery is as follows:
    
    Suppose $X/S$ is smooth. Since $\mathcal{P}(S)$ is cartesian closed, the functor $h_SX\times-: \mathcal{P}(S)\to \mathcal{P}(S)$ preserves colimits. By \Cref{dense prod}, the functor $h_SX\times-$ preserves dense open immersions. It preserves $\mathbb{A}^1$-projections (by definition of $\mathbb{A}(S)$) and Nisnevich covers. Hence it preserves $bir_n(S)$. Since $\mathcal{P}(S)$ is generated under colimits by smooth $X/S$ and products commute with colimits, we deduce that for any space $\esc{X}\in \mathcal{P}(S)$, the product functor $\esc{X}\times-: \mathcal{P}(S)\to \mathcal{P}(S)$ preserves $bir_n$. Hence, if $\esc{Y} $ is another space in $\mathcal{P}(S) $, the composite $$\esc{X}\times \esc{Y}\to\esc{X}\times L ^n\esc{Y} \to L^n\esc{X}\times L^n \esc{Y}$$ is an $n$-birational equivalence. But since $\mathcal{H}^{n}(S)$ is closed under limits, the codomain of the composite is an $n$-birational local space. Hence, $L^n\esc{X}\times L^n \esc{Y}$ is an $n$-birational localization of $\esc{X}\times\esc{Y}$, meaning that $L^n\big(\esc{X}\times\esc{Y}\big)\to L^n\esc{X}\times L^n \esc{Y}$ is an equivalence, which is to say that $L^n$ preserves products. 
    Suppose $X/S$ is smooth. Since $\mathcal{P}(S)$ is cartesian closed, the functor $h_SX\times-: \mathcal{P}(S)\to \mathcal{P}(S)$ preserves colimits. Now, by \Cref{dense prod}, the functor $h_SX\times-$ preserves dense open immersions. It preserves $\mathbb{A}^1$-projections (by definition of $\mathbb{A}(S)$) and Nisnevich covers. Hence it preserves $bir_n(S)$. Since $\mathcal{P}(S)$ is generated under colimits by smooth $X/S$ and products commute with colimits, we deduce that for any space $\esc{X}\in \mathcal{P}(S)$, the product functor $\esc{X}\times-: \mathcal{P}(S)\to \mathcal{P}(S)$ preserves $bir_n$. Hence, if $\esc{Y} $ is another space in $\mathcal{P}(S) $, the composite $$\esc{X}\times \esc{Y}\to\esc{X}\times L ^n\esc{Y} \to L^n\esc{X}\times L^n \esc{Y}$$ is an $n$-birational equivalence. But since $\mathcal{H}^{n}(S)$ is closed under limits, the codomain of the composite is an $n$-birational local space. Hence, $L^n\esc{X}\times L^n \esc{Y}$ is an $n$-birational localization of $\esc{X}\times\esc{Y}$, meaning that $L^n\big(\esc{X}\times\esc{Y}\big)\to L^n\esc{X}\times L^n \esc{Y}$ is an equivalence, which is to say that $L^n$ preserves products. 


    The arguments before and after this remark are intended to specify that we have an operadically correct localization.
 \end{rem}

The pointed version is similar:
\begin{thm}
    For every  scheme $S$, the chain of localization functors $$\mathcal{P}(S)_\bullet\xrightarrow[]{L_{mot}} \mathcal{H}_\bullet^{\mathbb{A}^1}(S)\xrightarrow[]{L^n} \mathcal{H}_\bullet^{n}(S)$$ can be lifted to a similar one of monoidal $\infty$-categories $$\mathcal{P}(S)_\bullet^{\wedge}\xrightarrow[]{L_{mot}} \mathcal{H}_\bullet^{\mathbb{A}^1, \wedge^{\mathbb{A}^1}}(S)\xrightarrow[]{L_{}^{n}} \mathcal{H}_\bullet^{n,\wedge^n} (S).$$ We thus have a commutative diagram of monoidal $\infty$-categories and monoidal functors:
    \[
    \xymatrix{
    \mathcal{P}(S)^{\times}\ar[r]^-{L_{mot}} \ar[d]&\mathcal{H}^{\mathbb{A}^1}(S)^\times\ar[r]^{L^n} \ar[d]&\mathcal{H}^{n}(S)^\times\ar[d]\\
    \mathcal{P}(S)_\bullet^{\wedge}\ar[r]^-{L_{mot
    }} &\mathcal{H}_\bullet^{\mathbb{A}^1, \wedge^{\mathbb{A}^1}}(S)\ar[r]^{L^n} &\mathcal{H}_\bullet^{n,\wedge^n} (S)\\
    }
    \]
\end{thm}
\begin{proof}
    Use [\cite{robalo2012noncommutative}, Corollaries 5.4, 5.6, and 5.9]. We just mention the definitions of the individual smash products. In $\mathcal{P}_\bullet(S)$, the smash product of $(\esc{X},x)$ and $(\esc{Y},y)$ is given by the cofiber in $\mathcal{P}_\bullet(S)$ of the canonical map from the coproduct to the product:
    $$(\esc{X},x)\vee (\esc{Y},y)\to (\esc{X},x)\times (\esc{Y},y)\to (\esc{X},x)\wedge (\esc{Y},y).$$
    Then one defines $\wedge^{\mathbb{A}^1}:=L_{mot}\circ \wedge$ and $\wedge^n:=L^n\circ \wedge$. 

    The vertical morphisms are similarly the derivations of the monoidal functors $(-)_\bullet$. 
\end{proof}The following is standard, and we state it without proof:
\begin{prop}
    The adjunction $\esc{X}\times-:\mathcal{H}^{n}(S)\rightleftarrows\mathcal{H}^{n}(S): (-)^{\esc{X}}$ descends to an adjunction $$\esc{X}\wedge^n-:\mathcal{H}^{n}_\bullet(S)\rightleftarrows\mathcal{H}^{n}_\bullet(S): (-)_\bullet^{\esc{X}}$$ where $(-)_\bullet^{\esc{X}}$ is the fiber of the map $(-)^{\esc{X}}\to (-)$ induced by the pointing $*\to \esc{X}$. In other words, the smash product monoidal structure on $\mathcal{H}^n_\bullet(S)$ is closed.
\end{prop}
\section{Schematic functorialities of \texorpdfstring{$\mathcal{H}^{n}$}{hb} }

We now turn to the main aspect of this paper, namely the analysis of the functoriality of the construction $S\mapsto \mathcal{H}^n(S)$ from $\mathrm{Sch}$ to $Pr^L$ for every $n\geq -1$. However, since $\mathcal{H}^{-1}(S)$ is the category $\mathcal{I}$ for every $S$, it suffices to study the functoriality for $n\geq 0$. When $n=0$, \S\ref{Hb=H0} shows that the functoriality of $\mathcal{H}^0$ yields that of $\mathcal{H}^b$ and vice versa.

Recall that given a morphism $f:S\to T$ of schemes, there is an adjunction $$f^*:\mathcal{P}(T)\leftrightarrows \mathcal{P}(S):f_*,$$ where the left adjoint is uniquely determined by the property that it preserves colimits and sends $h_T(X)$ to $h_S(X\times_TS)$. This makes the assignment $S\mapsto \mathcal{P}(S)$ into a presheaf $\mathcal{P}(-)^*:\mathrm{Sch}^{op}\to Pr^L$ (in fact, a presheaf with values in $\mathcal{T}op_\infty$, the $\infty$-category of large $\infty$-toposes) given by $f\mapsto f^*$.

Since any topology $\sigma$ is, by definition, stable under pullback, this presheaf induces another presheaf: $\mathcal{P}_{\sigma}(-)^*:\mathrm{Sch}^{op}\to Pr^L$ given by $S\mapsto \mathcal{P}_{\sigma}(S)$ and $f\mapsto L_\sigma f^*$. In particular, the Nisnevich topos construction defines a presheaf $\mathcal{P}_{nis}(-)^*:\mathrm{Sch}^{op}\to Pr^L$ (in fact, a presheaf taking values in $\mathcal{T}op_\infty$). 

On the other hand, when $f$ is smooth, precomposition by $f$ on objects of $Sm_S$ defines another functor $f_\shrp: \mathcal{P}(S)\to \mathcal{P}(T)$, which is the canonical left adjoint to $f^*$ and is uniquely determined by the fact that it sends $h_S(X)\mapsto h_T(X)$ and preserves colimits. Thus, the assignment $S\mapsto \mathcal{P}(S)$ also extends to a presheaf $\mathcal{P}(-)_\shrp:\mathrm{Sch}^{sm}\to Pr^L$ given by $p\mapsto p_\shrp$. It is trivial that $p_\shrp$ preserves Nisnevich coverings; hence we obtain another presheaf $\mathcal{P}_{nis}(-)_\shrp:\mathrm{Sch}^{sm}\to Pr^L$ given by $S\mapsto \mathcal{P}_{nis}(-)$ and $p\mapsto L_{nis}p_\shrp$. 

To summarize, let us note that there is a $(2,1)$-category of correspondences on $\mathrm{Sch}$, usually denoted $Corr(\mathrm{Sch})_{sm,all}$, with additional vertical morphisms given by smooth maps. The above discussions (with the standard exchange formulas) can be summarized as follows:
\begin{prop}
    There is a presheaf $$\mathcal{P}_{nis}(-):Corr(\mathrm{Sch})_{sm,all}^{op}\to Pr^L$$ whose restriction to horizontal maps recovers the presheaf $\mathcal{P}_{nis}(-)^*$, while its restriction to vertical maps recovers the presheaf $\mathcal{P}_{nis}(-)_\shrp$.
\end{prop}

The crux of the functoriality of the ordinary motivic formalism is that, because $\mathbb{A}^1$-projections are stable under both $f^*$ (for arbitrary $f$) and $p_\shrp$ (for all smooth $p$), the above presheaf localizes further to a presheaf taking values $S\mapsto \mathcal{H}^{\mathbb{A}^1}(S)$:
\begin{prop}
 The assignment $S\mapsto \mathcal{H}^{\mathbb{A}^1}(S)$ extends to a presheaf $$\mathcal{H}^{\mathbb{A}^1}(-):Corr(\mathrm{Sch})_{sm,all}^{op}\to Pr^L$$ whose restriction to the horizontal maps recovers the presheaf $\mathcal{H}^{\mathbb{A}^1}(-)^*$ given by $f\mapsto L_{mot}f^*$, while the restriction to the vertical maps recovers the presheaf $\mathcal{H}^{\mathbb{A}^1}(-)_\shrp$ given by $p\mapsto L_{mot}p_\shrp$.
\end{prop}
In this section, we plan to conduct a similar study of the unstable $n$-birational motivic homotopy category construction $S\mapsto \mathcal{H}^n(S)$. Throughout this section, $n$ denotes a natural number $\geq 0$.
\subsection{Pullback along morphisms lifting generalizations}


Unlike the ordinary motivic case, the pullback functors $f^*$ in general do not preserve $n$-dense open immersions. Openness is not really the issue; the issue is density and the appropriate behavior of the height of the density. The following lemma is the best we have; we make a definition for this.
\begin{defn}\label[defn]{uglt defn}
We say that a morphism of schemes is universally generalization lifting, or UGLT for short, if its arbitrary base change lifts generalizations.
\end{defn}
\begin{ex}
Universally open maps (for example, flat maps) and universal homeomorphisms are the primary examples. The composition of UGLT morphisms is UGLT.
\end{ex}
It follows that the wide class of UGLT morphisms of $\mathrm{Sch}$ forms a subcategory $\mathrm{Sch}^{uglt}$. Given a Qcqs scheme $S$, we denote by $\mathrm{Uglt}_S$ the full subcategory of $\mathrm{Sch}_{/S}$ spanned by UGLT morphisms. Moreover, $\mathrm{Uglt}^{uglt}_S$ denotes the slice category $({\mathrm{Sch}^{uglt})}_{/S}$, i.e., the wide subcategory of $\mathrm{UGLT}$ schemes and $\mathrm{UGLT}$ morphisms over $S$. 

Note that, since $\mathrm{Sch}^{uglt}$ does not have a terminal object, it is not possible to recover $\mathrm{Sch}^{uglt}$ as one of its slices. Therefore, studying presheaves on $\mathrm{Uglt}_{S}$ for an arbitrary base scheme is not the most general setup one might expect (in contrast to the ordinary motivic case). Therefore, we will state some of the results (the absolute ones) on $\mathrm{Sch}^{uglt}$, while some (the relative ones) will be stated over $\mathrm{Uglt}_S$ for an arbitrary base scheme $S$.

The following is the key to the announced functoriality of the $n$-birational motivic construction.
 \begin{lem}\label[lem]{glt pull}
    Suppose $f: S\to T$ lifts generalizations. Let $U\hookrightarrow T$ be an $n$-dense open immersion. Then $f^{-1}U\hookrightarrow S$ is also $n$-dense.
\end{lem}

\begin{proof}
We may safely assume that $U$ is an open subscheme of $T$ and may as well choose $f^{-1}U$ to be an open subscheme of $S$. We have to show that $f^{-1}U\subset S$ is $n$-dense. To see this, let $s\in S-f^{-1}(U)$. This implies that $f(s)\in T-U$ and thus has codimension at least $n+1$. Any maximal chain of generalizations starting at $f(s)$ must therefore have length at least $n+1$. Since $f$ lifts generalizations, we may lift one such maximal chain to obtain a chain of generalizations starting at $s$ and thus having length at least $n+1$. This implies that $cod_S(S-f^{-1}(U))\geq n+1$. 
\end{proof}
\begin{lem}\label[lem]{ugltpull}
     Let $f: S\to T$ be a universally generalization-lifting morphism of schemes. Then pullback along $f$ preserves $n$-dense open immersions. 
\end{lem}
\begin{proof}
       Let $U\to X\in B_n(T)$; we have to show that $f^{-1}(U\to X)\in B_n(S)$. By base change to $X$, we reduce to the previous lemma and the glt morphism $X\times _TS\to X$. 
\end{proof}

We thus have the following pullback functoriality:
\begin{prop}\label[prop]{generic push}
    Let $f: S\to T$ be a morphism that universally lifts generalizations. Then the pushforward along $f$ preserves $n$-birational motivic spaces, and, dually, the pullback functor $f^*$ preserves $n$-birational motivic equivalences. Thus, the adjunction below on the left descends to the one on the right: $$f^*:\mathcal{P}(T)\leftrightarrows\mathcal{P}(S):f_*\text{,  }L^nf^{*}:\mathcal{H}^{n}(T)\leftrightarrows \mathcal{H}^{n}(S): f_*.$$
\end{prop}
\begin{proof}
   Since $f^*$ is a left adjoint and $bir_n(T)$ is the saturation class of $\mathbb{A}(S)\cup B_n(S)$ and $nis(T)$, it suffices to show that it takes $\mathbb{A}(T)\cup B_n(T)$ to $\mathbb{A}(S)\cup B_n(S)$ and $nis(T)$ to $nis(S)$. The preservation of the set of $\mathbb{A}^1$-projections is obvious from its definition, whereas \Cref{ugltpull} settles the case of $B_n$. Since Nisnevich covers are stable under pullbacks, preservation of $nis$ is formal. Consequently, $f^*$ preserves $n$-birational motivic equivalences (for all $n$). Therefore, the right adjoint $f_*$ preserves $n$-local objects, and hence the result.
\end{proof}


\begin{cor}\label[cor]{flat pull push}
    If $f:S\to T$ is flat, then $f^*$ preserves $n$-birational equivalences, and we have an induced adjunction $L^nf^{*}:\mathcal{H}^{n}(T)\leftrightarrows \mathcal{H}^{n}(S): f_*$.
\end{cor}
\begin{cor}\label[cor]{stalks of n bir equiv}
    Let $\esc{X}\to \esc{Y}$ be a morphism of $S$-spaces that is a $n$-birational motivic equivalence over $S$. Then for every point $s\in S$, the morphism $\esc{X}_{{\mathcal{O}_{S,s}}}\to \esc{Y}_{{\mathcal{O}_{S,s}}}$ (resp. its base change to $\mathcal{O}_{S,s}^h$ or $\mathcal{O}_{S,s}^{sh}$) is a $n$-birational equivalence in $\mathcal{P}(\mathcal{O}_{S,s})$ (resp. $\mathcal{P}(\mathcal{O}_{S,s}^h)$ or $\mathcal{P}(\mathcal{O}_{S,s}^{sh})$).
\end{cor} 
\begin{proof}
    This follows from \Cref{flat pull push} using the fact that $Spec(\mathcal{O}_{S,s})\to S$ (resp. $Spec(\mathcal{O}_{S,s}^h)\to S$ or $Spec(\mathcal{O}_{S,s}^{sh})\to S$) is a flat morphism of schemes (being a limit of etale maps; see the paragraph above \Cref{continuity of Hb}).
\end{proof} 
\begin{proof}
    Flat maps lift generalizations [\cite[\href{https://stacks.math.columbia.edu/tag/03HV}{Tag 03HV}]{stacks-project}]. Thus, the statement follows from the above remark.
\end{proof}
\begin{prop}
    If $f:S\to S' $ is a morphism of schemes that universally lifts generalizations, then the induced functor $f_*:\mathcal{H}^n(S)\to \mathcal{H}^n(S')$ preserves all filtered colimits. When $n=0$, this moreover preserves sifted colimits. 
\end{prop}
\begin{proof}
    Since both categories are generated under colimits by compact (projective when $n=0$) objects [\Cref{compact generation}], namely schemes smooth over the base $S'$, and since $f^*$ takes these generators to schemes smooth over the base $S$, and thus to compact (projective when $n=0$) objects in particular. The claim then follows by adjunction.
\end{proof}

The case of smooth extension is obviously straightforward:
\begin{prop}\label[prop]{smooth extension pull}
    If $p:R\to S$ is a smooth morphism of schemes, the smooth extension map $L_{mot }p_{\shrp} : \mathcal{H}^{\mathbb{A}^1}(R)\to \mathcal{H}^{\mathbb{A}^1}(S) $ (or, $p_{\shrp} : \mathcal{P}(Sm_R)\to \mathcal{P}(Sm_S) $) preserves height $n$-birational equivalences. Thus, smooth pullback commutes with $n$-birational localizations: $L^np^*\simeq p^*L^n$, and we have adjunctions 
   \[ \begin{array}{ccc}
  L^np_{\sharp} :\mathcal{H}^{n}(R)\leftrightarrows \mathcal{H}^{n}(S): p^* & &p^*:\mathcal{H}^n(S)\leftrightarrows \mathcal{H}^n(R):p_*
\end{array}
\]
\end{prop}
\begin{proof}
    Being a dense open immersion or an $\mathbb{A}^1$ projection is a property independent of the structure maps of the associated schemes. Thus, $p_{\shrp} (B_n(R))\subset  B_n(S)$ and $p_{\shrp} (\mathbb{A}(R))\subset  \mathbb{A}(S)$. Moreover, $p_\shrp$ trivially preserves all nisnevich coverings. Since $p_\shrp$ preserves all colimits, it maps the saturation class $bir_n(R)$ of $\mathbb{A}(S)\cup B_n(R)$ and $nis(R)$ into the saturation class $bir_n(S)$ of $\mathbb{A}(S)\cup B_n(S)$ and $nis(S)$. Since $p^*$ is the right adjoint, it follows that, dually, it preserves $n$-birational motivic spaces. The second adjunction is a formal consequence. 
    
    
    Since, by \Cref{generic push}, $p^*$ also preserves n-birational local equivalences, we know that $p^*\esc{X}\to p^*L^n\esc{X}$ is an $n$-birational equivalence for every $\esc{X}\in \mathcal{P}(S)$. But from the above paragraph, we know that $p^*L^n\esc{X}$ is $n$-birational. Therefore, $p^*L^n\esc{X}$ is the $n$-birational localization of $p^*\esc{X}$, and so $L^np^*\esc{X}\simeq p^*L^n\esc{X}$. 

    The last adjunction is just the one from \Cref{pull adjunction}, because $$L^np^*_{|\mathcal{H}^n(S)}\simeq p^*L^n_{|\mathcal{H}^n(S)}\simeq p^*_{|\mathcal{H}^n}(S).$$
    \end{proof}
    \begin{rem}
Note that for $n=0$, by \Cref{bir local is P1 local}, in all the proofs above, after tracing the dense equivalences, it suffices to trace only the Nisnevich equivalences. Moreover, over qcqs schemes, by \Cref{nis is bir}, this can be further reduced to tracing the $\bigsqcup$-equivalences. In fact, it is easier to trace $\bigsqcup$-equivalences than Nisnevich ones.
\end{rem}
To use these observations fluently, we make some definitions.
\begin{defn}
    Given a functor $F: \mathcal{P}(T)\to \mathcal{P}(S)$ between the presheaf toposes of schemes $S,T$, define the $n$-Birational motivic derivation of $F$ as $$F(n):=L^n_SFb^n_T:\mathcal{H}^n(T)\to \mathcal{H}^n(S),$$ where $L^n_S$ is the $n$-birational motivic localization functor and $b_T^n$ denotes the inclusion $\mathcal{H}^n(T)\subset \mathcal{P}_{nis}(T).$
    
   In particular, for a morphism $f:S\to T$ of schemes, define the morphism: \begin{flalign*}
       f^{n*}:=f^*(n)\equiv L^n_Sf^* b_T^n: \mathcal{H}^n(T)\to \mathcal{H}^n(S)\\ f_{*n}:=f_*(n)=L^n_Tf_*b^n_S: \mathcal{H}^n(S)\to \mathcal{H}^n(T).
       \end{flalign*}
       If $f$ is smooth, we use the notation $$f_{n\shrp}:=f_\shrp(n)\equiv L^n_Tf_{\shrp} b^n_S.$$
\end{defn}
\begin{rem}
When $f$ is UGLT, \Cref{generic push} implies $f_{n*}\simeq {f_*}_{|\mathcal{H}^n(S)}$. On the other hand, when $f$ is smooth, \Cref{smooth extension pull} shows that $f^{n*}\simeq {f^*}_{|\mathcal{H}^n(T)}$.
\end{rem}
\begin{rem}\label[rem]{pull adjunction}
    From \Cref{ugltpull}, whenever $f$ is UGLT, $f^{n*}$ has a right adjoint given by $f_{n*}\simeq {f_*}_{|\mathcal{H}^n(S)}$, and the push-pull adjunction is now expressed as $f^{n*}:\mathcal{H}^{n}(T)\leftrightarrows \mathcal{H}^{n}(S): f_*$. Since $h_T(X)\to h^n_T(X)$ is an $n$-birational equivalence by construction, the same theorem implies that $f^*h_T(X)\to f^*h^n_T(X)$ is an $n$-birational equivalence. Hence, $$L^nf^*h^n_T(X)\simeq L^nf^*h_T(X)\simeq L^nh_S(X\times_TS).$$ That is, $f^{n*}h^n_T(X)\simeq  h^n_T(X\times_T S)$. The existence of a right adjoint and the presentability of $\mathcal{H}^n(T)$ imply that this left adjoint $f^{n*}$ is uniquely determined by the assignment $h^n_T(X/T)\mapsto h^n_T(X\times_T S)$ and by preservation of colimits. 
\end{rem}
\begin{rem}\label[rem]{extension adjunction}
    Similarly, when $f$ is smooth, $f_{n\shrp}$ is a left adjoint to $f^{n*}\simeq {f^*}_{|\mathcal{H}^n(T)}$. This left adjoint is uniquely determined by the fact that it preserves colimits and maps $h^n_S(W)$ to $h^n_T(W)$ for every smooth $W/S$. 
\end{rem}
  \begin{thm}[{Smooth base change}]\label{smooth base change}
    Suppose we have a pullback square of schemes    \[
    \xymatrix{
    V\ar[r]^{ q'}\ar[d]_{p'}&Y\ar[d]^{P}\\
    U \ar[r]_{q}&X
    }
    \]
  where the vertical maps are smooth and the horizontal ones are UGLT. Then the exchange transformations of $n$ birational motivic functors are equivalences:
  \begin{flalign}
      p'_{n\shrp} q'^{n*}\to q^{n*}p_{n{\shrp} }  \text{      from } \mathcal{H}^n(Y) \text{ to } \mathcal{H}^n(U)\\
      p^{n*}q_{n*}\to q'_{n*}p'^{n*} \text{      from } \mathcal{H}^n(U) \text{ to } \mathcal{H}^n(Y)
  \end{flalign}
   
    \end{thm}
    \begin{proof}
The second transformation is the adjoint mate of the first, so it suffices to show that the natural transformation in (2) is an equivalence.

Since $p_{\shrp}'$ preserves $n$-birational equivalences (see \Cref{smooth extension pull}), it follows that $p_{\shrp}'q^{'*}\to p'_\shrp L^nq^{'*}$ is an $n$-birational equivalence. Hence $L^n(p_{\shrp}'q^{'*})\simeq L^np'_\shrp L^nq^{'*}$. By our definition, the left-hand side of (2) is precisely $L^np'_{\shrp} L^nq'^*$. Thus, we identify the left-hand side with $L^n(p_{\shrp}'q^{'*})$. Similarly, since $q^*$ preserves $n$-birational equivalences by \Cref{generic push}, the right-hand side of (2) is $L^nq^*L^np_{\shrp} \simeq L^nq^*p_{\shrp}$. Therefore, the transformation in (2) is simply $$L^n\big({p'_\shrp q'^{*}}_{|\mathcal{H}^n(Y)}\to {q^*p_\shrp}_{|\mathcal{H}^n(Y)}\big).$$ It is thus sufficient to verify that $p'_\shrp q^{'*}\to q^*p_\shrp$ is a natural equivalence of functors $\mathcal{P}(Sm_Y)\to \mathcal{P}(Sm_U)$. Since $\mathcal{P}(Sm_Y)$ is generated by $h_Y(W)$'s for smooth $W/Y$ and all the functors involved in the last natural transformation ($p'_\shrp q^{'*}\to q^*p_\shrp$) are left adjoints, it suffices to verify that the components of this last transformation on these $h_Y(W)$ are equivalences. But clearly, $$p'_\shrp q^{'*}(h_Y(W))\cong h_U(V\times_YW),\text{ while }q^*p_\shrp h_Y(W)= q^*h_X(W)\cong h_U(U\times_XW).$$ So the claim follows by the pasting lemma for pullbacks.
    \end{proof}
    \begin{cor}\label[cor]{fully faithful j_*}
    For an open immersion $U\xhookrightarrow{j} X$, the birational pushforward ${j_{*}}_{|\mathcal{H}^n(U)}: \mathcal{H}^n(U)\to \mathcal{H}^n(X)$ and the birational extension $j_{n\shrp}$ are fully faithful.
    \end{cor}
    \begin{proof}
        Since open immersions are monomorphisms, we have a pullback square that satisfies the conditions of the theorem above:   
        \[
    \xymatrix{
    U\ar[r]^{1}\ar[d]_{1}&U\ar[d]^{j}\\
    U \ar[r]_{j}&X
    }
    \]
    Thus, the natural transformation from the above theorem [\Cref{smooth base change} (2)], namely $1\to j^{n*}j_{n\shrp} $ (which, by definition, is the unit of the adjunction), is an equivalence. This implies that the left adjoint $j_{n\shrp} $ (of $j^{n*}\simeq j^*_{|\mathcal{H}^n(U)}$) is fully faithful. Then, as a formal consequence, ${j_{*}}_{|\mathcal{H}^n(U)}$, being a right adjoint of $j^*_{_{|\mathcal{H}^n(U)}}$, is also fully faithful.
\end{proof}
    

    The first difference between the $\mathbb{A}^1$-homotopical world and the birational world appears in the following theorem. It's a simple observation that the $\mathbb{A}^1$-homotopy category is not $\mathbb{A}^1$-local (in the sense that the pullback along the projection $p:\mathbb{A}^1_S\to S$ induces an equivalence). For example, the pullback of $h_{\mathbb{A}_k^1}\mathbb{G}_{m,k}$ along the zero section $0\to \mathbb{A}^1_k$ is empty, i.e., the initial object. But $h_{\mathbb{A}^1_k}\emptyset$ cannot be motivically equivalent to $h_{\mathbb{A}^1}\mathbb{G}_{m,k}$ over $\mathbb{A}^1_k$. Indeed, $\mathbb{G}_{m,k}$ is $\mathbb{A}^1$-rigid in $\mathcal{P}(\mathbb{A}^1_k)$ while being a nonempty sheaf of sets. 
  \vspace{.4cm}  
    \begin{rem}\label[rem]{motivic category is a1 faithful}
        However, the pullback map $p^*:\mathcal{H}^{\mathbb{A}^1}(S)\to\mathcal{H}^{\mathbb{A}^1}(\mathbb{A}^1_S)$ is fully faithful. Indeed, the map $p_{\shrp} p^*\to 1$ is a motivic equivalence. To see this, note that both functors preserve colimits, so it suffices to check this on smooth $X/S$, where it is simply the projection $\mathbb{A}^1_S\times_SX\to X$.
    \end{rem}

    \begin{thm}\label{n dense equiv}
   For an $n$-dense open immersion $j:V\xhookrightarrow{}S$, the smooth extension map $$j_{i\shrp} :\mathcal{H}^{i}(V)\to \mathcal{H}^{i}(S)$$ is an equivalence for all $i\leq n$. Consequently, $j_{i*}\simeq {j_*}_{|\mathcal{H}^i(V)}$ and $j^{i*}\simeq j^*_{|\mathcal{H}^i(S)}$ are also equivalences. 
    \end{thm}
        \begin{proof}
       By \Cref{fully faithful j_*}, it suffices to show that the natural transformation $\varepsilon: j_{i\shrp } j^{i*}\to 1$ of endofunctors of $\mathcal{H}^i(S)$ is an equivalence. Since both functors preserve colimits and $\mathcal{H}^i(S)$ is generated under colimits by $h^i_S(W)$ for smooth $p:W\to S$, it is enough to show that the transformation is an equivalence on these generators. Now, 
       \begin{flalign}
           j_{i\shrp}j^{i*}h^i_S(W)&=L^ij_\shrp (h^i_V(W_V)) \text{ [by \Cref{pull adjunction}]}\\&=L^ih_S(W_V) \text{ [by \Cref{extension adjunction}]}
       \end{flalign}
       Clearly, the map $\varepsilon_{h^i_S(W)}$ is given by $L^ih_S(W_V)\to L^ih_S(W)$, which can be identified as $L^i (j':X_V\to X )$, where $j'$ is the pullback of $j$ along $p$. Since $p$ is smooth and thus lifts generalizations, and $j$ is $i$-dense for all $i\leq n$, by \Cref{ugltpull} we have that $j'$ is an $i$-dense open immersion and thus $L^i (j':X_V\to X )$ is indeed an equivalence as required.
    \end{proof}
\begin{cor}[Dense locality of $\mathcal{H}^0$]\label[cor]{dense equiv}
For a dense open immersion $j:V\xhookrightarrow{}S$, the smooth extension map $$j_{0\shrp}:=L_{0}j_\shrp :\mathcal{H}^0(V)\to \mathcal{H}^0(S)$$ is an equivalence, and hence so are $j_*\simeq L_{0}j_*$ and $j^*\simeq L_{0}j^*$. A similar statement holds for $\mathcal{H}^b:=L_{dense}\mathcal{P}_\Sigma$. 
\end{cor}
\begin{proof}
    The first statement follows immediately from \Cref{n dense equiv}. The last claim follows from \Cref{Motivic equivalences are Birational equivalences}, which identifies $\mathcal{H}^0\equiv \mathcal{H}^b$.
    \end{proof} 
\begin{rem}\label[rem]{Ldense is dense inv}
    Following the proof, it is easy to see that the statement above holds without motivic localization. That is, $$L_{dense}j_\shrp :L_{dense}\mathcal{P}(V)\to L_{dense}\mathcal{P}(S)$$ is an equivalence of $\infty$-categories whenever $j$ is dense. 
\end{rem}
Here is an interesting corollary:

\begin{cor}
    Suppose $f:\esc{X}\to \esc{Y}$ is a morphism of $S$-spaces. Then $f$ is an $n$-birational motivic equivalence if and only if there exists an $n$-dense open $V\subset S$ such that $f_V$ is an $n$-birational motivic equivalence in $\mathcal{P}(V)$. 
\end{cor} 

\begin{cor}
    Suppose $f:X\to Y$ is a morphism of smooth schemes over $S$ that has $n$-birationally contractible fiber over an $n$-dense open $V\subset Y$, i.e., $f_V:X_V\to V$ is an $n$-birational motivic equivalence in $\mathcal{P}(V)$. Then $f$ is an $n$-birational equivalence in $\mathcal{P}(S)$.
\end{cor}
\begin{proof}
    From the corollary above, it is clear that $f$ is a $0$-birational motivic equivalence in $\mathcal{P}(Y)$. Since the structure map $p: Y\to S$ is smooth and $f\in \mathcal{P}(S)$ is the image $p_\shrp f$ of $f\in \mathcal{P}(Y)$ under $p_\shrp:\mathcal{P}(Y)\to \mathcal{P}(S)$, the claim follows from \Cref{smooth extension pull}.
\end{proof}
\textbf{Presheaf description of $\mathcal{H}^{n}(-)$}

To summarize this section, note that for a composition $g\circ f$ with $f$ a UGLT morphism, we have $f^{n*}\circ g^{n*}\simeq (g\circ f)^{n*}$. Indeed, since $f^*$ preserves $n$-birational equivalences by \Cref{generic push}, the map $f^*g^*\to f^*L^ng^*$ is an $n$-birational equivalence, so that $$L^n(f^*g^*)\simeq L^n(f^*L^ng^*)\simeq L^nf^*\circ L^ng^*.$$ Moreover, when $g$ is also UGLT, so that both $f^{n*}$ and $g^{n*}$ are left adjoints, the equivalence $f^{n*}\circ g^{n*}\simeq (g\circ f)^{n*}$ holds, in fact, in $Pr^L$. Therefore, we may define:
    \begin{defn}
The presheaf $\mathcal{H}^{n}(-)^*: \mathrm{Sch}^{uglt,op}_{}\to Pr^L$ is given by the assignment $S\mapsto \mathcal{H}^{n}(S)$ and $f\mapsto f^{n*}$. 
\end{defn}
\begin{lem}\label[lem]{^* is cart}
The presheaf $$\mathcal{H}^{n}(-)^*: \mathrm{Sch}^{uglt, op}\to Pr^L$$ extends to a presheaf $$\mathcal{H}^{n}(-)^{\times,*}: \mathrm{Sch}^{uglt,op}\to \mathrm{CAlg}(Pr^{L,\otimes}).$$ In other words, for any UGLT morphism of schemes $f:S\to T$, the functor $f^{n*}: \mathcal{H}^n(T)\to \mathcal{H}^n(S)$ is (strictly) monoidal with respect to the cartesian monoidal structure on $\mathcal{H}^n(S)$ and $\mathcal{H}^n(T)$.
\end{lem}
\begin{proof}
Since $f^{n*}\equiv L^n_Sf^*b^n_T$, it suffices to show that each component is Cartesian. Because $b_T^n$ is a right adjoint, it is clearly Cartesian. On the other hand, $f^*$, being left exact, is also Cartesian. Finally, $L^n_S$ is Cartesian by \Cref{Ln is cart}.
\end{proof}
\begin{cor}\label{cart pull functoriality}
    If $T$ is a scheme, then $\mathcal{H}^{n} (-)^{\times,*}$ induce a presheaf of $\mathcal{H}^n(T)^\times $-modules given by $$\mathcal{H}^{n} (-)^{\times,*}:\mathrm{Uglt}_T^{uglt,op}\to \mathrm{Mod}_{\mathcal{H}^{n}(T)^{\times}}(Pr^{L,\otimes}).$$
\end{cor}
\begin{defn}
    Similarly, define the co-presheaf $$\mathcal{H}^n(-)_\shrp: \mathrm{Sch}^{sm}\to Pr^L$$ by $S\mapsto \mathcal{H}^{n}(S)$ and $p\mapsto p_{n\shrp}$.
\end{defn}
\begin{rem}
We must caution the reader that this is not a co-presheaf of cartesian monoidal categories. In other words, $p_{n\shrp}$ need not be a cartesian functor. In fact, since its right adjoint $p^{n*}$ is cartesian, $p_{n\shrp}$ is only op-lax monoidal [\cite{lurie2017higher}, Corollary 7.3.2.7]. This op-lax monoidal structure is monoidal for open immersions, as the next proposition shows.
\begin{prop}\label[prop]{j_sharp is cart}
Let $j: U\hookrightarrow X$ be an open immersion. Then the embedding $j_{n{\shrp}}:\mathcal{H}^n(U)\hookrightarrow\mathcal{H}^n(X)$ lifts to a monoidal embedding $j_{n{\shrp}}:\mathcal{H}^n(U)^\times \hookrightarrow\mathcal{H}^n(X)^\times $.
\end{prop}
\begin{proof}
The claim follows similarly to [\cite{khan2016motivic}, Proposition 6.1.4]. We shall sketch the argument below. By [\cite{lurie2017higher}, Corollary 2.4.1.8], it suffices to show that $j_{n\shrp}$ preserves finite products. Since $L^n$ is cartesian, it suffices to show this for $j_\shrp$ at the level of presheaves. But $j_\shrp$ is a left adjoint and hence preserves colimits, and the cartesian product functor preserves colimits in each variable. By the compact generation of $\mathcal{P}(U)$ by smooth schemes, it suffices to prove this on a pair $(h_UX,h_UY)$ for smooth schemes $X,Y/U$. This is clear since the diagonal of $j$ is an isomorphism.
\end{proof}
\end{rem}


It follows immediately from \Cref{^* is cart} that:

\begin{prop}[Smooth projection formula]\label{sm proj formula}
Let $p: Y\to X$ be a smooth morphism of schemes, and let $\esc{X}\in \mathcal{H}^{n}(X)$ and $\esc{Y}\in \mathcal{H}^{n}({Y})$. Then the canonical maps \begin{flalign}p_{\shrp}(p^*\esc{X}\times\esc{Y})\to \esc{X}\times p_{\shrp}\esc{Y}\\
{(p_\shrp\esc{Y})}^{\esc{X}}\to p_\shrp(\esc{Y}^{p^*\esc{X}})\end{flalign} are equivalences in $\mathcal{H}^{n}(X)$ and $\mathcal{H}^{n}(Y)$, respectively. 
\end{prop}
\begin{proof}
By adjunction, the equivalences are dual to each other, so it is enough to verify the first. But the first follows from the cartesian-ness of $L^n$ [\Cref{Ln is cart}] and the smooth projection formulas at the level of presheaves. Again, use the fact that each functor involved in the formulas preserves colimits and that the formula holds on the generators of $\mathcal{P}(X)\times \mathcal{P}(Y)$, namely $(h_X(W),h_Y(W'))$ for smooth $W/X$ and $W'/Y$. 
\end{proof}
\begin{cor}\label{cart extension functoriality}
    Let $X$ be a scheme, and consider the category $\mathrm{Uglt}_X^{sm}$ of schemes UGLT over $X$ with morphisms the smooth ones. The assignment $p\mapsto p_{n\shrp}$ makes $\mathcal{H}^{n}_{\shrp}(-)$ into a $\mathcal{H}^n(X)^{\times}$-module co-presheaf. Formally, $$\mathcal{H}^{n}(-)_\shrp: \mathrm{Uglt}_X\to \mathrm{Mod}_{\mathcal{H}^n(X)^\times }(Pr^{L,\otimes}).$$
\end{cor}
\begin{proof}
    Let $r$ be a smooth morphism of schemes UGLT over $X$ defined by the following diagram:
    \[
    \xymatrix{
    Y\ar[rr]^r\ar[dr]_p&&Z\ar[dl]^q\\
    &X&
    }
    \]
Assume that $\esc{X}\in \mathcal{H}^n(X)$ and $\esc{Y}\in \mathcal{H}^n(Y)$. We must show that $r_{n\shrp}(p^{n*}\esc{X}\times \esc{Y})\simeq q^{n*}\esc{X}\times r_{n\shrp}\esc{Y}$. This follows immediately from the above proposition. Indeed, $$r_{n\shrp}(p^{n*}\esc{X}\times \esc{Y})\simeq r_{n\shrp}(r^{n*}q^{n*}\esc{X}\times \esc{Y})\simeq q^{n*}\esc{X}\times r_{n\shrp}\esc{Y}$$ (the first equivalence holds because $p^{n*}\simeq r^{n*}\circ q^{n*}$, and the last follows from \Cref{sm proj formula}).
\end{proof}

It is possible to write similar statements for the pointed setups. We leave the details to the reader. See [\cite{khan2016motivic}] for a similar treatment. 

The content of this section can then be summarized as follows. Let $\mathrm{Corr}(\mathrm{Sch}_S)_{ sm,uglt}$ be the $(2,1)$-category whose objects are schemes, vertical morphisms are smooth maps, horizontal morphisms are uglt, and composition of morphisms is given by composition through pullback completion of spans. Because a pullback requires a choice up to an isomorphism, this is forced to be a $(2,1)$-category.

\begin{rem}
Note that `birational morphisms' in the original sense of algebraic geometry are actual morphisms in this category, with both the vertical and horizontal morphisms being dense open immersions.
\end{rem}

\begin{thm}\label{full functoriality}
    There is  a presheaf $$\mathcal{H}^n(-): \mathrm{Corr}
    (\mathrm{Sch})_{sm,uglt}^{op}\to Pr^L$$ whose restriction to the vertical maps recovers $\mathcal{H}^n(-)_\shrp: \mathrm{Sch}^{sm}\to Pr^L$ while the restriction to the horizontal maps recovers $\mathcal{H}^{n}(-)^*: \mathrm{Sch}^{uglt,op}_{}\to Pr^L$.
\end{thm}
\begin{proof}
This follows from \Cref{smooth base change} and [\cite{MR3701352}, Theorem 2.1.3].
\end{proof}
\begin{cor}
Let $S$ be a Qcqs scheme. Then the presheaf $$\mathcal{H}^n(-):\mathrm{Corr}(\mathrm{Sch})_{sm,uglt}^{op}\to Pr^L$$ extends to 
 a presheaf $$\mathcal{H}^n(-)/\mathcal{H}^n(S):Corr(\mathrm{Uglt}_S)_{sm,uglt}^{op}\to Mod_{\mathcal{H}^n(S)^\times}(Pr^{L,\otimes})$$ whose restriction to vertical maps recovers $$\mathcal{H}^n(-)_\shrp: \mathrm{Uglt}_S^{sm}\to Mod_{\mathcal{H}^n(S)^\times}(Pr^{L,\otimes})$$ from \Cref{cart pull functoriality}, while the restriction to horizontal maps recovers $$\mathcal{H}^{n}(-)^*: \mathrm{Uglt}^{uglt,op}_{S}\to Mod_{\mathcal{H}^n(S)^\times}(Pr^{L,\otimes})$$ from \Cref{cart extension functoriality}.
\end{cor}


\begin{rem}
       It is possible to write down all the results of this subsection for the naive Birational Homotopy construction $\mathcal{H}^b(-):=L_{dense}L_\Sigma\mathcal{P}(-)$, recovering the functor $\mathcal{H}^{0}$ over Qcqs schemes. And in fact, just as above, this is an easier task than the (however, equivalent) Birational \textbf{Motivic} counterpart $\mathcal{H}^{0}$.
\end{rem}
\subsection{Refined pushforward}\label{3.2}
As we have mentioned earlier, $n$-dense open immersions need not be stable under arbitrary pullbacks. The problem is easily visible when pulling back along closed immersions: it might happen that the pullback of a dense immersion becomes empty. For example, if $i: Z\rightarrowtail X$ is a closed immersion of a divisor, then it fails to pull back the dense open subscheme $j:U=X\setminus Z\hookrightarrow X$. Of course, when the closed immersions are UGLT, the results of the previous section do apply:
\begin{cor}
    Let $i:Z\to S$ be a UGLT closed immersion. Then $i^*$ preserves $n$-birational motivic equivalences, and $i_*$ preserves $n$-birational motivic spaces.
\end{cor}


In the example above, the observation is made in the wrong light. The correct observation is that the pullback along $i$ takes $0$-dense open immersions to $(-1)$-dense open immersions. The following result is meant to indicate that, in general, instead of asking for stability of dense open immersions under closed pullback, the right question is to ask what height of density a closed immersion takes to what height after pulling back. For the rest of this section, we assume that $S$ is Noetherian Universally Catenary (NUC for short). We denote the subcategory of $\mathrm{Sch}$ consisting of NUC schemes as $\mathrm{Nuc}$. 
  \begin{thm}\label{closed rel d push}
Suppose $Z\xhookrightarrow{i} S$ is a closed immersion of $NUC$ schemes of maximal codimension bounded by $d$, i.e., $mcod(i)\leq d$. Recall that this means $i$ maps generic points to points of codimension at most $d$. Then for every $n\geq 0$, the functor $i^*:\mathcal{P}(S)\to \mathcal{P}(Z)$ takes $(d+n)$-birational motivic equivalences to $n$-birational motivic equivalences, and dually, $i_*$ takes $n$-birational motivic spaces to $(n+d)$-birational motivic spaces. We have an adjunction $$L_{n}i^*: \mathcal{H}^{{n+d}}(S) \leftrightarrows \mathcal{H}^{n}(Z):i_{*}.$$
    \end{thm}
        \begin{proof}
       Since $i^*$ preserves Motivic equivalences, it suffices to show that it maps $(d+n)$-dense open immersions to $n$-dense open immersions.
        
         Let $f:X\to S$ be smooth, and $U\hookrightarrow X$ be an $(d+n)$-dense open immersion. We must show that the open immersion $$(Z\times_SX)\bigcap U\hookrightarrow Z\times_SX$$ is $n$-dense. 
         
         Note that the inclusion $Z\times_SX\hookrightarrow X$ also has maximum codimension bounded by $d$. Indeed, if $\tilde{z}\in (Z\times_SX)^{(0)}$, then, since every scheme is locally noetherian and $f$ is flat, we have $$cod_X(\tilde{z})=cod_S(f(\tilde{z}))+cod_{X_{{f(\tilde{z})}}}(\tilde{z})$$ [\cite{bachmann2019voevodsky}, Theorem 2.1]. Since $\tilde{z}\in X_{{f(\tilde{z})}}$ ($\subset Z\times_SX$) is a generic point of $Z\times_SX$, it is also a generic point of $X_{{f(\tilde{z})}}$. Thus the last term, $cod_{X_{{f(\tilde{z})}}}(\tilde{z})$, is zero, and we have $cod_X(\tilde{z})=cod_S(f(\tilde{z}))$. But $f_Z: Z\times_SX\to Z$ being flat implies that $f(\tilde{z})$ is a generic point of $Z$. Since the maximum codimension of $Z$ in $S$ is $\leq d$, it follows that $cod_S(f(\tilde{z}))\leq d$ and thus $cod_X(\tilde{z})\leq d$.
          
          
Moreover, since $X\to S$ is of finite type and $S$ is universally catenary, we know that $X$ is catenary. By the above paragraph and base change, we may thus reduce to showing that $(d+n)$-dense open immersions into the base $S$ pull back to $n$-dense open immersions into $Z$. Since every $\eta_Z\in Z^{(0)}$ has codimension at most $d$, for any $(d+n)$-dense open subscheme $U\hookrightarrow S$, we have $i(Z^{(0)})\subset U$. The pullback $Z\bigcap U\hookrightarrow Z$ therefore satisfies $Z^{(0)}\subset Z\bigcap U$, i.e., $Z\bigcap U$ is dense in $Z$. Now let $z\in Z$ be a point with $cod_Z(z)\leq n$. Let $z\leftsquigarrow \eta_z$ be a generalization with generic point $\eta_z\in Z^{(0)}$. Since $S$ is catenary, by [\cite[\href{https://stacks.math.columbia.edu/tag/02I6}{Lemma 02I6}]{stacks-project}] we have $$cod_Sz=cod_S\eta_z+cod_{\bar{{\eta_{z}}}} z\leq cod_S\eta_z+cod_{Z} z\leq  d+n.$$  Because $U\hookrightarrow S$ is $(d+n)$-dense, it follows that $z\in U$. Hence $z\in U\bigcap Z$. Since $z\in Z$ was arbitrary with $cod_Z(z)\leq n$, this observation implies that $Z\bigcap U\hookrightarrow Z$ is $n$-dense.
          
          Thus, $i^*$ takes $(d+n)$-birational motivic equivalences to $n$-birational motivic equivalences, and dually, $i_*$ takes $n$-birational motivic spaces to $(d+n)$-birational motivic spaces. As a formal consequence we deduce the adjunction $$L^{n}i^*: \mathcal{H}^{{n+d}}(S) \leftrightarrows \mathcal{H}^{n}(Z):{i^{}_{*}}_{|\mathcal{H}^{n}(Z)}.$$
\end{proof}
\begin{cor}
    Let $f:\esc{X}\to S$ be an $n$-birational motivic equivalence over an NUC scheme $S$. Then for every $s\in S^{(n)}$, the fiber $\esc{X}_{k(s)}$ is birationally contractible over $k(s)$.
\end{cor} 
\begin{proof}
    Let $i_s:\bar{x}\to S$ be the closed immersion of the closure of $s$ in $S$. Then, by \Cref{closed rel d push}, $i_{s}^*f$ is a $0$-birational motivic equivalence over $\bar{x}$. The claim now follows from \Cref{stalks of n bir equiv}.
\end{proof}
\begin{thm}\label{closed push full-faithful}
   Let $Z\xhookrightarrow{i} S$ be a closed immersion of NUC schemes. Then the functor $$i_*:  \mathcal{H}^{n}(Z) \to \mathcal{H}^{{n+mcod(i)}}(S)$$ induced by \Cref{closed rel d push} is fully faithful.
    \end{thm}
        \begin{proof}
      Clearly, ${i_*}_{|\mathcal{H}^{n}(Z)}={i^{mot}_*}_{|\mathcal{H}^{n}}$.
We know that $i_*^{mot}$ is fully faithful [Corollary 7.4.3, \cite{khan2016motivic}]. So we win.    \end{proof}

For a closed immersion $i$, let us define $$i^{*}_n:=L_{{n-mcod(i)}}i^*:\mathcal{H}^{n}(S)\to \mathcal{H}^{n-mcod(i)}(S).$$ 
\begin{rem}
However, because $mcod$ is not additive, it is not possible to arrange closed immersions with their $mcod$ grading in a suitable category. Hence, this result alone cannot be used to extend the construction $S\mapsto\mathcal{H}^-(S)$ to a (say, mcod-graded) presheaf, as we have done for uglt maps. This is possible when $mcod$ is additive, for example, when $mcod=cod$, i.e., for equidimensional closed immersions.
\end{rem}
\begin{cor}
    If $\delta(i)=0$, then $i^*$ maps $(n+cod(i))$-birational equivalences to $n$-birational equivalences.
\end{cor}
Since codimension is additive for catenary schemes [\cite[\href{https://stacks.math.columbia.edu/tag/02I6}{Lemma 02I6}]{stacks-project}], by the above corollary, the following definition gives a presheaf $$\mathcal{H}^\mathbb{N}(-)^*: \mathrm{Nuc}^{closed_{\delta=0},op}\to Cat_\infty$$ via the following assignment:
\begin{flalign*}
    &S\mapsto \underset{n\in \mathbb{N}}{\bigsqcup }\mathcal{H}^{n}(S) \\
    &i \mapsto (n\mapsto i^{*}_n)
\end{flalign*}
    

\section{Continuity properties of \texorpdfstring{$\mathcal{H}^{n}$}{hb}}

Let us go back to the presheaf $\mathcal{P}_{nis}(-)^*: \mathrm{Sch}^{op}\to Pr^L$ consisting of presheaves satisfying Nisnevich descent. It turns out that this presheaf has good geometric properties, most of which descend to the motivic homotopy presheaf $\mathcal{H}^{\mathbb{A}^1}(-)^*:Sch^{op}\to Pr^L$. 

For example,

\begin{enumerate}
    \item $\mathcal{P}_{nis}(-)^{*}$ itself satisfies Nisnevich descent. With a little work, one can show that this holds even for $\mathcal{H}^{\mathbb{A}^1}(-)^*$; see [\cite{hoyois2017six}, Proposition 4.8] and [\cite{khan2016motivic}, Proposition 6.1.6] (below we will briefly recall how it works). 
    \item The presheaves $$\mathcal{P}(-)^*, \mathcal{P}_{nis}(-)^*, \mathcal{H}^{\mathbb{A}^1}(-)^*:Sch^{aff,op}\to Pr^L$$ take filtered colimits (i.e., pro schemes in $\mathrm{Sch}$ with affine transition maps) to filtered colimits (in $Pr^L$). In short, they are affine-continuous.
\end{enumerate}
    In fact, $\mathcal{H}^{\mathbb{A}^1}(-)^*$ has a few more properties that do not hold without the $\mathbb{A}^1$-localization:
\begin{enumerate}[resume]    
    \item $\mathcal{H}^{\mathbb{A}^1}(-)^*:Sch^{op}\to Pr^L$ is a deformation invariant.
    \item $ \mathcal{H}^{\mathbb{A}^1}(-)^*:Sch^{op}\to Pr^L$  takes $\mathbb{A}^1$ projections to fully faithful functors.
\end{enumerate}


In this section, we will briefly recall each property and then study analogous properties of the $n$-Birational motivic homotopy presheaf: $$\mathcal{H}^n(-)^{*}:\mathrm{Sch}^{uglt,op}\to Pr^L.$$
\subsection{Nisnevich descent for $\mathcal{H}^{n}(-)^*$}
We begin our discussion of the continuity properties of the $n$-birational motivic construction by showing that even this candidate satisfies Nisnevich descent. This is neither surprising nor unnecessary, given that the Nisnevich topos construction itself has the same property.  
\begin{prop}[Nisnevich descent of $\mathcal{H}^n(-)^*$]\label{bir detect} 

       \begin{enumerate}
           \item For a Nisnevich cover $\{f_\alpha: U_\alpha\to S\}$, the pullback maps $f_{\alpha}^*: \mathcal{P}(Sm_S)\to \mathcal{P}(Sm_{U_\alpha})$ preserve and detect $n$-birational motivic equivalences.
           \item  The presheaf $\mathcal{H}^{n}(-)^*:\mathrm{Sch}^{ uglt, op}\to Cat_\infty$ is Nisnevich separated, i.e., for a Nisnevich cover $f_\alpha: U_\alpha\to S$, the pullback maps $$f_\alpha^{n*}\simeq L^nf_\alpha^{*}: \mathcal{H}^{n}(S)\to \mathcal{H}^{n}(U_\alpha)$$ are jointly conservative. 
           \item The presheaf $\mathcal{H}^{n}(-)^{*}:\mathrm{Sch}^{ uglt, op}\to Pr^L$ from \Cref{full functoriality} is a Nisnevich sheaf.
       \end{enumerate}  
\end{prop}    
\begin{proof}
1. The only-if part of (1) follows from the fact that each $f_\alpha$ is etale and hence uglt. For the if part of (1), let $f:U=\bigsqcup U_\alpha \to S$ be the disjoint union of the cover. Consider the \v{C}ech nerve $c_\bullet f: \check{C}_\bullet f \to S$ and the associated \v{C}ech resolution $\varepsilon: C_\bullet\to 1$ of endo-functors of $\mathcal{P}(Sm_S)$. In other words, for every $[m]\in \Delta$, define $$C_m (F)= (c_mf)_{\shrp} (c_mf)^{*}(F)\in \mathcal{P}(Sm_S).$$ Clearly, for each $F\in \mathcal{P}(Sm_S)$, the map  $|\varepsilon_F|: \colim_{}C_\bullet F\to F$  is a Nisnevich equivalence, since this holds when evaluated on smooth schemes. 
      
Now let $h:F\to G$ be a morphism in $\mathcal{P}({Sm_S})$ such that for every $\alpha$, the morphism $f_{\alpha}^{*}h$ is an $n$-birational equivalence in $\mathcal{P}(Sm_{U_\alpha})$. Since $f^*h=\bigsqcup_\alpha f_\alpha^*h$, and $bir_n(U)$ is colimit-closed, we have $f^*h\in bir_n(U)$. Recall that for each $m$, $$c_mf=f\times _Sf\times_S\cdots \times_Sf\text{ (the $(m+1)$-fold self-pullback of }f\text{),}$$ so that $c_mf=f\phi$ for some etale map $\phi$. Because $n$-birational motivic equivalences are stable under pullbacks by etale maps [\Cref{generic push}] and $(c_mf)^*h=(f\circ \phi)^*h\simeq \phi^*(f^*h)$, it follows that $(c_mf)^*h$ is an $n$-Birational motivic equivalence in $\mathcal{P}(\check{C}_mf)$. 


Moreover, since $c_mf$ is smooth, \Cref{smooth extension pull} implies that $(c_mf)_\shrp$ preserves $n$-birational motivic equivalences. Hence, $(c_mf)_{\shrp} (c_mf)^{*}h$ is an $n$-birational motivic equivalence in $\mathcal{P}(S)$. Consequently, $C_\bullet h$ is a simplicial $n$-birational motivic equivalence. Now consider the commutative square below,        \[
    \xymatrix{
    \colim_{}C_\bullet F \ar[r]^-{\varepsilon_F}\ar[d]_{ \colim_{}C_\bullet h}&F\ar[d]^{h}\\
 \colim_{}C_\bullet G \ar[r]_-{\varepsilon_G}&G
    }
    \]
Because $bir_n(S)$ is closed under colimits, we deduce from the last paragraph that the left vertical map in the square above is an $n$-birational motivic equivalence. The first paragraph, on the other hand, guarantees that both horizontal maps in the square are Nisnevich equivalences, and hence $n$-birational motivic equivalences. By the 2-out-of-3 property of $bir_n(S)$, the commutativity of the square implies that the right vertical map, namely $h$, is also an $n$-birational motivic equivalence.

2. Statement (2) follows from the first one.

3. For (3), first note that, since $Pr^L\subset Cat_\infty$ is closed under limits [\cite{lurie2009higher}, Proposition 5.5.3.13], by a general theory of descent it suffices to show that, for a scheme $S$, the restriction of the presheaf $\mathcal{H}^n$ to $Et_S$ is a Nisnevich sheaf of $\infty$-categories. By the 2-out-of-three property of etale maps, all morphisms of $Et_S$ are etale, and hence, by \Cref{smooth extension pull}, for any $f:A\to B\in Mor(Et_S)$, $$\mathcal{H}^{n*}(f)\equiv L_{n}f^*\simeq f^*_{|\mathcal{H}^{n}}.$$ Thus, the restriction $\mathcal{H}^{n}(-)^*_{|Et_S}$ is a subpresheaf of the presheaf $\mathcal{P}_{nis}(-)_{|Et_S}$. Since $\mathcal{P}_{nis}(-)^*_{|Et_S}$ is a sheaf for the Nisnevich topology [\cite{hoyois2017six}, proposition 4.8 (1)], and the functor $$\mathcal{H}^{n}(-)^*_{|Et_S}\subset \mathcal{P}_{nis}(-)_{|Et_S}$$ is fully faithful, it suffices, by stability of fully faithful functors under limits, to show that, given a Nisnevich cover as in (1), the family $f_{\alpha}^{ *}$ can detect $n$-birational locality of Nisnevich sheaves of spaces. For this, we need to show that if $f_{\alpha}^{ *} \esc{F}$ is $n$-birational local for each $\alpha$, then so is the given Nisnevich sheaf $\esc{\esc{F}}$. This follows from the fact that the $\mathbb{A}^1$ projection morphisms, as well as $n$-dense open immersions [\Cref{ugltpull}], are stable under pullback by smooth maps. 

For example, suppose $V\xhookrightarrow{}X$ is an $n$-dense open immersion of smooth schemes over $S$. Let us show that $\esc{F}^X\to \esc{F}^V$ in $\mathcal{P}(Sm_S)$ is an equivalence. Since both $\esc{F}^X$ and $\esc{F}^V$ are Nisnevich local, it suffices to show that $\esc{F}^X\to \esc{F}^U$ is a Nisnevich equivalence. Since $\{f_\alpha:U_\alpha\to S\}$ is a Nisnevich cover of $S$, by [\cite{hoyois2017six}, Proposition (1)] it suffices to show that, for every $\alpha$, the map $f_\alpha^{*}(\esc{F}^X\to \esc{F}^V)$ is a Nisnevich local equivalence in $\mathcal{P}(U_\alpha)$. But $$f_\alpha^*(\esc{F}^X)\simeq (f_\alpha^*\esc{F})^{X\times_SU_\alpha}\text{ and similarly, }f_\alpha^*(\esc{F}^V)\simeq (f_\alpha^*\esc{F})^{V\times_SU_\alpha}.$$ So we have to show that the map $$(f_\alpha^*\esc{F})^{U_\alpha\times_SX}\to(f_\alpha^*\esc{F})^{U_\alpha\times_S V} $$ induced by $U_\alpha\times_SV\to U_\alpha\times_SX$ is an equivalence. Since $f_\alpha : U_\alpha\to S$ is UGLT, it follows from \Cref{ugltpull} that the map $U_\alpha\times_SV\to U_\alpha\times_SX$ over $U_\alpha$ is an $n$-dense open immersion. On the other hand, $f_\alpha^*\esc{F}$ is given to be $n$-birational motivic local. So this last requirement holds by definition of $n$-Birational motivic locality.
\end{proof}
       \begin{cor}[$\mathcal{H}^n(-)^*$ is additive]\label{zar add}
           Suppose we have a finite disjoint union of schemes $X=\bigsqcup X_\alpha$ (for example, $X$ is unibranch, and the $X_\alpha$'s are its irreducible connected components). Then $$\mathcal{H} ^{n}(X)\xrightarrow[]{\Pi i_\alpha^*}\underset{\alpha} {\prod}  \mathcal{H}^{n}(X_\alpha){}$$ is an equivalence of $\infty$-categories. 
       \end{cor}
\begin{rem}\label[rem]{Ldense not nis sheaf}
    The spirit of \Cref{Ldense is dense inv}, however, does not apply to either of the results above. More precisely, $L_{dense}\mathcal{P}(-)$ is not even a contractible sheaf, let alone a Nisnevich or a $\sqcup$-sheaf. Indeed, $L_{dense}\mathcal{P}(\emptyset)$ is equivalent to $Spc$, not the terminal category. 
\end{rem}
\subsection{Topological invariance}
Let us recall that the small Nisnevich topos construction is a topological invariant:
\begin{lem}\label[lem]{def inv of small nis topos}
    The presheaf $Shv(-_{nis})^*:Sch^{op}\to Pr^L$ given by $S\mapsto Shv(S_{nis})$ sends universal homeomorphisms to equivalences. In particular, it is a deformation-invariant.
\end{lem}
\begin{proof}
    After one knows that radical maps are precisely the universally injective maps [\cite[\href{https://stacks.math.columbia.edu/tag/01S4}{Lemma 01S4}]{stacks-project}], the list of conditions in [\cite{MR1106918}, E.3] is equivalent to the corresponding morphism being a universal homeomorphism [\cite[\href{https://stacks.math.columbia.edu/tag/04DF}{Lemma 04DF}]{stacks-project}]. So this is just a restatement of [\cite{MR1106918}, E.3]. 
\end{proof}
Unfortunately, this does not extend to the Large Nisnevich topos construction $\mathcal{P}_{nis}(-)^*$. For example, suppose we are dealing with the reduction morphism $i:k=S^{red}\to S$ where $S=Spec (k[x]/x^2)$, then the unit $1\to i_*i^*$ fails to be an equivalence. Indeed, $i_*i^*\mathbb{A}^1_S\simeq i_*\mathbb{A}^1_S$ and $(i_*\mathbb{A}^1_S)(S)=Sm_S(S^{red},\mathbb{A}^1_S)=\Gamma(S^{red})=k$ while $\mathbb{A}^1_S(S)=\Gamma(S)=k[x]/x^2$. In fact, once the affine line is contracted, this issue completely disappears. That is, the motivic homotopy category construction is well known to be deformation-invariant:

\begin{thm}[Motivic Deformation Invariance]\label{motivic deformation invariant}
 The presheaf $\mathcal{H}^{\mathbb{A}^1}(-)^*: \mathrm{Sch}^{op}\to Pr^L$ sends thickenings of schemes to equivalences. Specifically, if $Z\xhookrightarrow{i} S$ is a nil immersion of schemes, then the functor $${i_*}_{mot}:={i_*}_{|\mathcal{H}^{\mathbb{A}^1}(Z)}:  \mathcal{H}^{\mathbb{A}^1}(Z) \to \mathcal{H}^{\mathbb{A}^1}(S)$$ is an equivalence of $\infty$-categories. In particular, $\mathcal{H}^{n}(X^{red})\to \mathcal{H}^{n}(X)$ is an equivalence.
    \end{thm}
\begin{proof}
For a closed immersion $i:Z\hookrightarrow X$, there is a cofiber sequence in $Pr^L$ [\cite{hoyois2021localization}, Remark 6]:
$$ \mathcal{H}^{\mathbb{A}^1}(X\setminus Z)\xrightarrow{j_\shrp}\mathcal{H}^{\mathbb{A}^1}(X)\xrightarrow{i^*}\mathcal{H}^{\mathbb{A}^1}(Z)$$
When $i$ is a thickening, we have $X\setminus Z=\emptyset$. Thus $\mathcal{H}^{\mathbb{A}^1}(X\setminus Z)=\mathcal{H}^{\mathbb{A}^1}(\emptyset)\simeq pt$, the zero object in $Pr^L$. Therefore, the cofiber of $\mathcal{H}^{\mathbb{A}^1}(X)\xrightarrow{i^*}\mathcal{H}^{\mathbb{A}^1}(Z)$ must be an equivalence.
\end{proof}
It follows that this induces nil-invariance for $\mathcal{H}^n(-)$ as well. To this end, we will need an $n$-birationality detection lemma for thickenings:
\begin{lem}\label[lem]{n bir detection by thickening}
    Let $i:Z\to S$ be a thickening. Then $i_*:\mathcal{P}_{nis}(Z)\to \mathcal{P}_{nis}(S)$ detects $n$-dense locality.
\end{lem}
\begin{proof}
    Let $j: U\hookrightarrow X$ be an $n$-dense open immersion over $Z$. By [\cite{grothendieck1966elements}, Proposition 18.1.1], there exists a Zariski cover $X_\alpha$ of $X$ such that $X_\alpha \simeq X_\alpha'\times_SZ$ for some $X'_\alpha\in Sm_S$. Since $X_\alpha\to X_\alpha'$ is a homeomorphism, by [\cite{grothendieck1966elements}, 18.1.2] the $n$-dense open immersion $j_\alpha :U_\alpha=U\times_XX_\alpha\hookrightarrow X_\alpha$ is pulled back from a unique open immersion $j_\alpha':U_\alpha'\hookrightarrow X_\alpha'$. Moreover, because both $X_\alpha\to X_\alpha'$ and $U_\alpha\to U_\alpha'$ are homeomorphisms, it is clear that $U_\alpha'\hookrightarrow X_\alpha'$ is also $n$-dense. This works for all finite fiber products. Moreover, by \Cref{dense prod} we know that the \v{C}ech nerves of both $J:=\sqcup j_\alpha$ and $J':=\sqcup j_\alpha'$ are simplicial $n$-birational motivic equivalences. Also, $\colim_{}\check{C}_\bullet (J)\simeq j$. 

    Now, suppose $i_*\esc{X}$ is $n$-dense local. Then we have an equivalence of morphisms: $$\mathrm{Map}(j,\esc{X})\simeq \lim_m \mathrm{Map} (\check{C} _mJ,\esc{X})\simeq \lim_m \mathrm{Map}(i^*\check{C} _mJ', \esc{X})\simeq \lim_m \mathrm{Map}(\check{C}_mJ', i_*\esc{X}).$$ This last morphism is an equivalence when $i_*\esc{X}$ is $n$-dense local, because $\check{C}_m j_\alpha'$ is $n$-dense for every $m$. Thus, the first morphism $\mathrm{Map}(j,\esc{X})$ is an equivalence, implying the $n$-dense locality of $\esc{X}$.
    \end{proof}

From this, we obtain the result of this subsection, namely the deformation invariance of the $n$-birational motivic homotopy category construction.
\begin{thm}[Deformation Invariance]\label{deformation invariant}
The presheaf $$\mathcal{H}^{n}(-)^*: Sch^{uglt,op}\to Pr^L$$ sends thickenings of schemes to equivalences of $\infty$-categories. Specifically, if $Z\xhookrightarrow{i} S$ is a nil immersion, then $i_*:  \mathcal{H}^{n}(Z) \to \mathcal{H}^{n}(S)$ is an equivalence of $\infty$-categories. In particular, the canonical map $$\mathcal{H}^{n}(X^{red})\to \mathcal{H}^{n}(X)$$ given by pullback along the reduction $X^{red}\subset X$ is an equivalence.
    \end{thm}
\begin{proof}
Since thickenings are UGLT, $i_*$ preserves $n$-birational motivic spaces, and we have an adjunction 
$$L^ni^*:\mathcal{H}^n(Z)\leftrightarrows \mathcal{H}^n(S):i_*(n):= {i_*}_{|\mathcal{H}^n(S)}={i_{mot}^*}_{|\mathcal{H}^n{(S)}}$$
Since ${i_{mot}^*}$ is an equivalence by \Cref{motivic deformation invariant}, it follows that $i_*(n)$ is also fully faithful. To see that it is essentially surjective, first note that if $\esc{X}\in \mathcal{H}^{n}(Z)\subset \mathcal{H}^{\mathbb{A}^1}(Z)$, then by \Cref{motivic deformation invariant}, there exists a motivic space $\esc{Y}\in \mathcal{H}^{\mathbb{A}^1}(S)$ such that $i_*\esc{Y}\simeq \esc{X}$. But then, \Cref{n bir detection by thickening} tells us that $\esc{Y}$ must be $n$-birational local.
\end{proof}
\begin{rem}\label[rem]{Ldense not def inv}
    The above result is likely false for the naive dense localization $L_{dense}\mathcal{P}(-)$ (as is the case in \Cref{Ldense not nis sheaf}) (even though, naive dense localization factors through the naive $\mathbb{A}^1$-localization).
\end{rem}
\subsection{Essentially smooth continuity for $\mathcal{H}^n$}
Let us now recall the continuity property for $\mathcal{P}(-)^*,\mathcal{P}^*_{nis}(-):\mathrm{Sch}^{op}\to Pr^L$.
\begin{prop}
Suppose $S=\plim{}S_\alpha$ is a pro-scheme with affine transition maps $f_{\alpha\beta}$, then the morpshim  $$\phi=\dcolim_{\alpha} \phi_\alpha^*:  \dcolim_{\alpha}\mathcal{P}(S_\alpha)\to \mathcal{P}(S)$$ defined in $Pr^L$, induced by the left adjoints $\phi_\alpha^*$ associated to the legs $\phi_\alpha: X\to X_\alpha$, is an equivalence. Equivalently, the morphism $$\plim_{\alpha} \phi_{\alpha,*}:  \mathcal{P}(S)\to \plim_{\alpha}\mathcal{P}(S_\alpha)$$ defined in $Cat_\infty$ is an equivalence. Similarly, $$\dcolim_{\alpha} L_{nis}\phi_\alpha^*:  \dcolim_{\alpha}\mathcal{P}_{nis}(S_\alpha)\simeq  \mathcal{P}_{nis}(S)$$ in $Pr^L$ and equivalently, $$\plim_{\alpha} \phi_{\alpha,*}:  \mathcal{P}_{nis}(S)\simeq \plim_{\alpha}\mathcal{P}_{nis}(S_\alpha ).$$ In short, the presheaves $\mathcal{P}(-)^*, \mathcal{P}^*_{nis}(-): \mathrm{Sch}^{op}\to Pr^L$ are pro-affine continuous.
\end{prop}
\begin{proof}
  Since colimits in $Pr^L$ are computed as limits in $Pr^R$, which in turn are computed in $Cat_\infty$, the `equivalently' parts of each statement are formal. Thus, it suffices to show the second equivalences in each case.
 
 Let us first address the case of $\mathcal{P}(-)^*$. We will show the second equivalence $$\plim_{\alpha} \phi_{\alpha,*}:  \mathcal{P}(S)\to \plim_{\alpha}\mathcal{P}(S_\alpha)$$ in $Cat_\infty$. To do so, we shall use standard cofiltering arguments, as in [\cite{hoyois2015quadratic}, Proposition C.7 (3)].
        
 Recall that all our smooth morphisms are finitely presented [see \S1.2]. By [\href{https://stacks.math.columbia.edu/tag/01ZM}{Lemma 01ZM} (1)], given a smooth $X/S$, we may assume there is an $\alpha_0$ and a finitely presented $X_{\alpha_0}\to S_{\alpha_0}$ such that $X\to S$ is pulled back from $X_{\alpha_0}\to S_{\alpha_0}$ via $S\to S_{\alpha_0}$. Moreover, by [\href{https://stacks.math.columbia.edu/tag/0C0C}{Lemma 0C0C}] and [\href{https://stacks.math.columbia.edu/tag/01ZQ}{Lemma 01ZQ}], we may find $\alpha_1$ such that the smooth and separated morphism $X_{\alpha_1}\to S_{\alpha_1}$ pulls back to the smooth (separated by assumption) map $X\to S$ via $S\to S_{\alpha_1}$. Thus $X=\plim_{\alpha\geq \alpha_1}X_\alpha$. Therefore, the pullback maps $\phi_\alpha^*: Sm_{S_\alpha}\to Sm_S$ induce an essential surjection $\plim \phi_\alpha^*: \plim Sm_{S_\alpha}\to Sm_S$. The fact that this functor is fully faithful is well known (use [\cite[\href{https://stacks.math.columbia.edu/tag/01ZB}{Section 01ZB}]{stacks-project}] with cofinality arguments, assuming everything is smooth). So, $  Sm_S\simeq\plim Sm_{S_\alpha}$. Since $\mathcal{P}$ commutes with limits in $Cat_\infty$, we deduce that $  \mathcal{P}(S)\simeq\plim \mathcal{P}({S_\alpha})$. Indeed, 
 \begin{flalign}
\mathcal{P}(S):= {\mathrm{\mathrm{Cat}}}_\infty(Sm_S^{op},Spc)&\simeq {\mathrm{\mathrm{Cat}}}_\infty((\plim Sm_{S_\alpha})^{op},Spc)\\&\simeq {\mathrm{\mathrm{Cat}}}_\infty(\dcolim Sm_{S_\alpha}^{op},Spc)\\&\simeq \plim {\mathrm{\mathrm{Cat}}}_\infty(Sm_{S_\alpha}^{op},Spc)=\plim \mathcal{P}(S_\alpha).     
 \end{flalign}
 

Similarly, descent for etale maps [\cite[\href{https://stacks.math.columbia.edu/tag/07RP}{Lemma 07RP}]{stacks-project}], open immersions [\cite[\href{https://stacks.math.columbia.edu/tag/0EUU}{Lemma 0EUU}]{stacks-project}], surjective maps [\cite[\href{https://stacks.math.columbia.edu/tag/07RR}{Lemma 07RR}]{stacks-project}], and isomorphisms [\cite[\href{https://stacks.math.columbia.edu/tag/081E}{Lemma 081E}]{stacks-project}] implies that a Nisnevich square on $X$ lifts from a Nisnevich square on some $X_\alpha$, and thus the above equivalence restricts to an equivalence $\mathcal{P}_{nis}(S) \simeq \plim \mathcal{P}_{nis}({S_\alpha})$.    
\end{proof}
Since affine projections on $S$ are trivially pulled back from an affine projection on some component $S_\alpha$, the above continuity extends to the presheaf $\mathcal{H}^{\mathbb{A}^1}(-)^*:\mathrm{Sch}^{op}\to Pr^L$, i.e., the canonical morphism $$  \dcolim_{\alpha}\mathcal{H}^{\mathbb{A}^1}(S_\alpha)\to \mathcal{H}^{\mathbb{A}^1}(S)$$ in $Pr^L$ is an equivalence [\cite{hoyois2015quadratic}, Proposition C.2(3)]. That is, the motivic homotopy category satisfies essentially smooth continuity for pro-smooth morphisms whose transition maps are affine. Unfortunately, in the birational setting, such a statement is not expected, as general affine maps need not be generalization-lifting. Because arbitrary UGLT maps might not behave well under limits, such a continuity result cannot be expected with only UGLT-affine transition maps.

To state the birational result, let us first note that, for a ring $R$, since $Ring_R\to Mod_R$ creates filtered colimits (Rings being an algebraic theory) and the category of flat modules $FlatMod_R\subset Mod_R$ is closed under filtered colimits, the category $Flat_R$ of flat rings over $R$ is also closed under filtered colimits in $Ring_R$. It follows immediately that if $S=\lim{}S_\alpha$ is a pro-limit in $\mathrm{Sch}$ with flat transition maps $\phi_{\beta \alpha}:S_\beta\to S_\alpha$, then each map $\phi_\alpha: S\to S_\alpha$ is flat. Therefore, in this situation, for all $n$, there is a well-defined morphism $\phi^{n*}_\alpha:\mathcal{H}^n(S_\alpha)\to \mathcal{H}^n(S)$ in $Pr^L$.
 \begin{thm}[Flat-affine continuuity of $\mathcal{H}^n(-)^*$]\label{continuity of Hb}
     The functor $\mathcal{H}^{n}(-)^*: \mathrm{Sch}^{flat,op}\to Cat_\infty$ is pro-affine continuous; i.e., if $S=\plim{}S_\alpha$ (with legs $\phi_\alpha: S\to S_\alpha$) is a pro-scheme with flat affine transition maps, then the pullback assembly map in $Pr^L$ $$\phi=\dcolim_{\alpha} \phi_\alpha^{n*}:  \dcolim_{\alpha}\mathcal{H}^n(S_\alpha)\to \mathcal{H}^n(S),$$ induced by the left adjoints $\phi_\alpha^*$, is an equivalence of presentable infinity categories. Equivalently, $$\plim_{\alpha} \phi_{\alpha,n*}:  \mathcal{H}^n(S)\to \plim_{\alpha}\mathcal{H}^n(S_\alpha)$$ is an equivalence in ${\mathrm{Cat}}_\infty$ (and thus in $Pr^L$).
    \end{thm}
    \begin{proof}
As mentioned earlier, we have $\mathcal{H}^{\mathbb{A}^1}(S)\simeq  \plim_{\alpha}\mathcal{H}^{\mathbb{A}^1}(S_\alpha)$. It then suffices to show that $n$-dense open immersions descend along the system.

 Now, if $U\xhookrightarrow{j} X$ is a dense open immersion (quasicompact by assumption), then [\cite[\href{https://stacks.math.columbia.edu/tag/01Z4}{Lemma 01Z4}]{stacks-project}] yields a qc open $j_{\alpha_{0}}:U_{\alpha_0} \subset X_{\alpha_0}$ such that $U_{\alpha_0}, X_{\alpha_0}$ are smooth and separated over $S_{\alpha_0}$ and $j_{\alpha_0}$ pulls back to $j$. Since open immersions have constructible image, all our transition maps are flat, and $j$ is dominant, we deduce from [\cite{grothendieck1966elements}, Proposition (8.10.3)] that for some $\alpha_1\geq \alpha_0$, $j_{\alpha_1}$ is dominant and hence a dense open immersion. Thus the above equivalence passes to $\mathcal{H}^0(S) \simeq \plim \mathcal{H}^0({S_\alpha})$. 


 To see the general $n$-birational case, suppose $j$ is of height $n$. Let $Z$ be the complement of $jU$. Consider the closure $Z_1$ of the image of $Z=X\setminus jU$ under the flat map $\phi_{\alpha_1}: X\to X_{\alpha_1}$. Clearly, $Z_1\subset(X_{\alpha_1}\setminus U_{\alpha_1}) $, and hence $\phi_{\alpha_1}^{-1}Z_1=Z$. Thus, $$\phi_{\alpha_1}^{-1}(X_{\alpha_1}\setminus Z_1)=X\setminus Z=U.$$ Let $U_{\alpha_1}':=X_{\alpha_1}\setminus Z_1$. Clearly, $U_{\alpha_1}\subset U_{\alpha_1}'$. By the denseness of $U_{\alpha_1}$ in $X_{\alpha_1}$, it follows that $ U_{\alpha_1}'$ is also dense in $X_{\alpha_1}$. By replacing $j_{\alpha_1}: U_{\alpha_1}\hookrightarrow X_{\alpha_1}$ with the open immersion $$j_{\alpha_1}':U_{\alpha_1}'=X_{\alpha_1}\setminus Z_1\hookrightarrow X_{\alpha_1},$$ we see that $j_{\alpha}'$ is a dense open immersion over $S_{\alpha_1}$ that pulls back to $j$ and such that $\phi_{\alpha_1}$ is surjective on the generic points of the closed complements. We claim that $j'_{\alpha_1}$ is $n$-dense. Indeed, by the methods of the second paragraph of \Cref{closed rel d push}, applied to the flat map $\phi_{\alpha_1}$, the codimension of the complement $Z_1=X_{\alpha_1}\setminus U_{\alpha_1}'$ satisfies $$cod_{X_\alpha}Z_1\geq cod_X(Z)\geq n+1.$$
 
 We have thus shown that, if $j\in B_{n}(S)$, then there exists $j_{\alpha_1}'\in B_n(S)$ that pulls back to $j$. From this discussion, the equivalence $$\mathcal{H}^{\mathbb{A}^1}(S)\simeq \plim_{\alpha}\mathcal{H}^{\mathbb{A}^1}(S_\alpha)$$ descends to an equivalence $$\mathcal{H}^n(S)\simeq \plim_{\alpha}\mathcal{H}^n(S_\alpha)$$ in ${\mathrm{Cat}}_\infty$, as desired.
\end{proof}
 This gives us an interesting consequence for the ($0$-)birational motivic homotopy category:
      \begin{cor}[Generic decomposition of $\mathcal{H}^0(-)$]\label{Hb of a var is Hb of its generic point}
         Let $X^{(0)}$ be the set of generic points of a qcqs unibranch scheme $X$. Then the canonical map $$\mathcal{H}^0(X)\to\underset{\eta\in X^{(0)}}{\prod}\mathcal{H}^0(k(\eta)),$$ induced by pullback along the pro-smooth map $\underset{\eta\in X^{(0)}}{\bigsqcup} \mathrm{Spec}(k(\eta))\hookrightarrow X$, is an equivalence of $\infty$-categories. In particular, for a reduced irreducible scheme $X$, we have an equivalence $\mathcal{H}^0(X)\simeq \mathcal{H}^0(K(X))$.
     \end{cor}
     \begin{proof}
First, by deformation invariance \Cref{deformation invariant}, we may assume that $X$ is reduced. Next, since $X$ is qcqs, it has finitely many connected components. Moreover, because $X$ is unibranch, each connected component is irreducible. Thus, by appeal to Zariski additivity [\Cref{zar add}], we reduce to the case of a reduced and irreducible scheme $V$.

In this case, we have the pro system $$SpecK(X)=\plim_{SpecA=U\in Dense_X}U$$ of dense open affine schemes of $X$. Because each $W$ in the limit is affine, the transition maps are indeed affine. Moreover, the transition maps, being open immersions, are in particular flat. The pro system $$D_X:=\{SpecA=U\in Dense_X\}$$ thus satisfies all the required conditions of the above theorem. Thus, from \Cref{continuity of Hb}, it follows that the map $$ \mathcal{H}^0(k(X))\to \plim_{\alpha}\mathcal{H}^0(U)$$ corresponding to the theorem, is an equivalence. But each transition map $\phi_{\alpha \beta}: U_\alpha\to U_\beta$ in $D_X$ as well as the inclusion $j_\alpha :U_\alpha \subset X$ is a dense open inclusion. So, by \Cref{dense equiv}, all the transition maps in $\plim_{\alpha}\mathcal{H}^0(U)$, namely $\phi_{\alpha \beta,*}:\mathcal{H}^0(U_\alpha)\to \mathcal{H}^0(U_\beta)$, are equivalences and under $j_\alpha^*$'s, each of these is equivalent to the identity functor of $\mathcal{H}^0(X)$. Consequently, $\plim_{\alpha}\mathcal{H}^0(U)\simeq \mathcal{H}^0(X)$, and we are done.
\end{proof}
\begin{cor}\label[cor]{Hb of a variety}
    Let $V$ be a variety over $k$. Then there is a canonical equivalence $\mathcal{H}^b(V)\to \mathcal{H}^b(K(V))$, where $\mathcal{H}^b := L_{dense}\mathcal{P}_\Sigma$ is the birational homotopy category construction of [\cite{0bat}] (also recalled in \S2.3 of this paper).
\end{cor}
\begin{proof}
We know from \Cref{Motivic equivalences are Birational equivalences} that $\mathcal{H}^b(-)=\mathcal{H}^0(-)$.
\end{proof} 
\begin{rem}\label[rem]{Ldense cont for var}
    As in \Cref{Ldense is dense inv}, the proof shows that the above equivalence holds even without motivic localization. More precisely, $L_{dnese}\mathcal{P}(V)\simeq L_{dense}\mathcal{P}(K(V))$. However, since $L_{dense}\mathcal{P}(-)$ is not a $\sqcup$-sheaf (see \Cref{Ldense not nis sheaf}) (and is likely neither a deformation-invariant, \Cref{Ldense not def inv}), the general case of \Cref{Hb of a var is Hb of its generic point} does not apply to $L_{dense}\mathcal{P}(-)$.
\end{rem}
\begin{cor}\label[cor]{fibersise criteria}
    Suppose $S$ is a Qcqs unibranch scheme and $f:\esc{X}\to \esc{Y}$ is a morphism of $S$-spaces. Then $f$ is a $0$-birational equivalence in $\mathcal{P}(S)$ if and only if its generic fibers are $0$-birational equivalences, i.e., for every generic point $\eta\in S^{(0)}$ the morphism $f_\eta:\esc{X}_{k(\eta)}\to \esc{Y}_{k(\eta)}$ is a $0$-birational equivalence in $\mathcal{P}(k(\eta))$.
\end{cor} 
\begin{cor}\label[cor]{fibersise criteria for smoo}
    Suppose $S$ is a Qcqs geometrically unibranch scheme and $f:X\to Y$ is a morphism in $Sm_S$ such that the generic fibers of $f$ are birationally contractible, i.e., for every generic point $\eta\in Y^{(0)}$ the scheme $f_\eta:X_{k(\eta)}\to k(\eta)$ is contractible in $\mathcal{H}^0(k(\eta))$. Then $f$ is a birational equivalence in $\mathcal{P}(S)$.
\end{cor} 
\begin{proof}
    From the corollary above, it is clear that $f$ is a birational equivalence in $\mathcal{P}(Y)$. If $q:Y\to S$ is the smooth structure map of $Y$, then $f\in \mathcal{P}(S)$ is the image of $f\in \mathcal{P}(Y)$ under $q_\shrp$. The claim thus follows from \Cref{smooth extension pull}.
\end{proof}
\subsection{Rationality property of \texorpdfstring{$\mathcal{H}^0(-)$}{h}}
In this subsection, we plan to show that, in the special case $n=0$, the corresponding birational homotopy category $\mathcal{H}^0$ (which is the birational homotopy category constructed in [\cite{0bat}]) is a birational presheaf of $\infty$-categories. Moreover, we show that this is homotopy-invariant and stably-rational-invariant, in the sense that it induces a fully faithful embedding along the appropriate projection maps.
 \begin{thm}\label{H^b is bir local}
 $\mathcal{H}^{0}(-)^{*}:\mathrm{Sch}^{uglt,op}\to Pr^L$ is a birational local presheaf, i.e., $$\mathcal{H}^{0}(-)^*\in L_{dense}\mathcal{P}_{\Sigma}(\mathrm{Sch}^{uglt}, Pr^L).$$
    \end{thm}
\begin{proof}
From \Cref{bir detect}, it is clear that $\mathcal{H}^{0}(-)^{*}$ lies in $\mathcal{P}_{\Sigma}(\mathrm{Sch}^{uglt}, Pr^L)$. The dense open invariance of $\mathcal{H}^{0}(-)^{*}$ follows from \Cref{dense equiv}.
\end{proof}
\begin{rem}
    Because of \Cref{dense equiv} and \Cref{nis is bir}, it is possible to give an alternative proof that $\mathcal{H}^{0}(-)^*: \mathrm{Sch}^{uglt,op}\to \mathrm{Cat}_\infty$ satisfies Nisnevich descent, provided one explicitly proves \Cref{zar add}, that is, that $\mathcal{H}^{0}(-)^*$ belongs to $\mathcal{P}_{\Sigma}(\mathrm{Sch}^{uglt}, \mathrm{Cat}_\infty)$.
\end{rem}
 
\begin{prop}[Homotopy invariance]\label{homotopy}
    For a scheme $S$, the pullback maps $p^{n*}: \mathcal{H}^{n}(S)\to \mathcal{H}^{n}(\mathbb{A}^m_S)$ induced by the projections $p:\mathbb{A}^m_S\to S$ are fully faithful.
\end{prop}
\begin{proof}We need to show that $p_{n\shrp}p^{n*}\to id$ is an equivalence. Since both are left adjoints, and hence preserve colimits, by \Cref{compact generation} it suffices to prove this on $h^n_S(X)$. But $p_{n\shrp}p^{n*}\to id$ evaluated on $h^n_S(X)$ is the projection map $L^n(h_S\mathbb{A}^1_X\to h_SX)$, which, by construction, is an equivalence in $\mathcal{H}^n(S)$. 
\end{proof}

\begin{rem}\label[rem]{rem about strong a1 inv}
The above proposition holds even for the $\mathbb{A}^1$ homotopy category $\mathcal{H}^{\mathbb{A}^1} (-)$; see \Cref{motivic category is a1 faithful}. However, as explained in the paragraph above the referenced remark, the $\mathbb{A}^1$ homotopy category $\mathcal{H}^{\mathbb{A}^1}(-)$ is not $\mathbb{A}^1$-invariant in the strongest sense, in that it does not send $\mathbb{A}^1$ projections to equivalences of $\infty$-categories. Initially, while developing the paper, the author expected that, for $\mathcal{H}^0(-)^*$, such a strong version of $\mathbb{A}^1$-locality is possible. The reason was \Cref{H^b is bir local} and the crucial observation of \Cref{bir local is P1 local} about birational presheaves of sets. At the end of this section, we will discuss a problem in combining these two results. We are unable to draw a final conclusion about the validity of such a strong version of $\mathbb{A}^1$-invariance of $\mathcal{H}^0(-)^*$. If this is true, then $\mathcal{H}^0(-)$ would be an example of a Motivic $\infty$-category in the sense of [\cite{kato2017motivic}]. 
\end{rem}
\begin{cor}[Rational invariance]\label{rational invariance}
Let $S$ be a scheme, and assume $q:X\to S$ is a rational scheme over $S$ (equivalently, $q$ is stably birational). Then the functor $q^{0*}=L^0q^*:\mathcal{H}^0(S)\rightarrow \mathcal{H}^0(X)$ is fully faithful.
    \end{cor}
    \begin{proof}
 By definition, there is a rational correspondence $X\xhookleftarrow{k} U\xhookrightarrow{j} \mathbb{A}^n_S$. Use \Cref{dense equiv} to obtain the composite equivalence: $$\mathcal{H}^0(X)\underset\simeq{\xleftarrow[]{k_{0\shrp}}}\mathcal{H}^0(U)\underset\simeq{\xleftarrow{j^*_b}} \mathcal{H}^0(\mathbb{A}^n_S).$$ It is not difficult to see that the morphism $q^{0*}$ in question is equivalent to the composite $$\mathcal{H}^0(X)\underset\simeq{\xleftarrow[]{k_{0\shrp}}}\mathcal{H}^0(U)\underset\simeq{\xleftarrow{j^*_b}} \mathcal{H}^0(\mathbb{A}^n_S)\xhookleftarrow{p^{0*}}\mathcal{H}^0(S)$$ which is fully faithful because, by \Cref{homotopy}, the pullback $\mathcal{H}^0(S)\xrightarrow{p^{0*}}\mathcal{H}^0(\mathbb{A}^n_S)$ is fully faithful.
\end{proof}
\begin{cor}
    The functor $\mathcal{H}^b(S)\to \mathcal{H}^b(\mathbb{P}^1_S)$ induced by the pullback along the projection map $\mathbb{P}^1_S\to S$ is fully faithful.
\end{cor}

\begin{cor}[Purely transcendence invariance]\label{trans}
For any field $k$, the purely transcendental embedding $k\subset k(t_1,\cdots,t_n)$ induces an embedding $\mathcal{H}^0(k)\subset \mathcal{H}^0(k(t_1,\cdots,t_n))$.
\end{cor}
\begin{proof}
    By continuity, \Cref{continuity of Hb} yields $\mathcal{H}^0(k(t_1,\cdots,t_n))\simeq \mathcal{H}^0(\mathbb{A}^n_k)$. The result then follows from \Cref{homotopy}.
\end{proof}
\begin{rem}
   Using \Cref{Ldense is dense inv}, \Cref{bir local is P1 local}, and \Cref{Ldense cont for var}, it is easy to observe that the overall method of proof of all the results above applies to $L_{dense}\mathcal{P}(-)$, implying its rational and pure transcendence invariance.
\end{rem}
We end this section with a discussion of what was already hinted at in \Cref{rem about strong a1 inv}. We recommend that the reader review that remark and return to this discussion. We plan to explain why the contents of \Cref{bir local is P1 local} and \Cref{dense equiv}\textbf{ cannot be used} to conclude that "the functor $p^{0*}:\mathcal{H}^0(S)\to \mathcal{H}^0(\mathbb{A}^1_S)$ is an equivalence of $\infty$-categories".

There are two major problems with combining these two mechanisms. The primary one is formal: the argument of \Cref{bir local is P1 local} applies to ($\mathcal{G}$-)sets (and thus even to $\infty$-groupoids), but not to ($\infty$-)categories. However, \Cref{homotopy} provides an immediate resolution, reducing the problem to a presheaf of large sets. That is, using \Cref{homotopy}, the trick to proving the claimed stronger $\mathbb{A}^1$-locality of $\mathcal{H}^0(-)^*$ should be to show that $\pi_0\mathcal{H}^0(S)\to \pi_0\mathcal{H}^0(\mathbb{A}^1_S)$ is surjective. This brings us to the second problem: the diagram that yields the required surjectivity has morphisms that do not lift generalizations. Therefore, the method of the second paragraph cannot be applied to the functor $\mathcal{H}^0: \mathrm{Sch}^{uglt,op}\to Cat_{\infty}$ alone. Here is the correct way to chase the diagram, though inconclusive. We proceed with the proof of \Cref{bir local is P1 local} and consider the following diagram (which makes sense since $i$, $j$, and $k$ are at least codimension $1$):

\[ \xymatrix{
  & \mathcal{H}^0(E) & & \\
  \mathcal{H}^0(S) \ar[ur]^{L^0\pi^*} & \mathcal{H}^1(W) \ar[d]^{L^1y^*}\ar[u]^{L^0j^*} & \mathcal{H}^0(L) \ar[ul]_{{\sim}^*} &\mathcal{H}^0(\mathbb{P}^1)\ar[l]^{{\sim}}  \\
  \mathcal{H}^1(\mathbb{P}^2) \ar[u]^{L^0i^*} \ar[ur]^{L^1a^*} \ar[r]^{L^1x^*}&\mathcal{H}^1(U) & \mathcal{H}^1(\mathbb{P}^1 \times \mathbb{P}^1) \ar[ul]_{L^1b^*} \ar[u]^{L_0k^*} & \mathcal{H}^1(\mathbb{P}^1)\ar[u]_{L^{0}_{\mathbb{P}^1}} \ar[ul]_{} \ar[l]_-{L^1p_1^*}
}\]
 Clearly, $k^*p_1^*\simeq id$. Since $k$ is of pure codimension 1, \Cref{closed rel d push} implies that $k^*$ takes $1$-birational motivic equivalences to $0$-birational motivic equivalences. Thus, $k^*p_1^*\to k^*L^1p_1^*$ is a $0$-birational motivic equivalence. Consequently, $$(L^{0}k^*)\circ (L^1p_1^*)=L^0(k^*L^1p_1^*)\simeq L^{0}(k^*p_1^*)\simeq L^{0}_{\mathbb{P}^1}.$$ But $L_{\mathbb{P}^1}^0$ is essentially surjective (being a localization). So $L^0k^*$ is also essentially surjective. On the other hand, since $k^*\simeq j^*b^*$ and $j$ is also of pure codimension $1$, we similarly have $$L^0k^*\simeq L^0(j^*b^*)\simeq L^0(j^*( L^1b^*))\simeq (L^0j^*)\circ (L^1b^*).$$ Thus, the essential surjectivity of $L^0k^*$ guarantees that $L^0j^*$ is also essentially surjective.

Again, since $j^*$ takes $1$-birational equivalences to $0$-birational equivalences, we get $$L^0(j^*a^*)\simeq  L^0(j^*(L^1a^*))\simeq (L^0j^*)\circ (L^1a^*).$$ Since $\pi^*i^*=j^*a^*$, this yields $$L^0(\pi^*\circ i^*)\simeq (L^0j^*)\circ (L^1a^*).$$ But since $\pi$ is UGLT, $\pi^*$ preserves birational equivalences by \Cref{generic push}. So, similarly $$L^0(\pi^*\circ i^*)\simeq L^0( \pi^*(L^0i^*))\simeq (L^0\pi^*)\circ (L^0i^*).$$ Combining the last two equivalences, we obtain $$(L^0j^*)\circ (L^1a^*)\simeq (L^0\pi^*)\circ (L^0i^*).$$ But in the last paragraph, we have seen that $L^0j^*$ is essentially surjective; so, to show that $L^0\pi^*$ is essentially surjective, it suffices, using this last equivalence, to show that $L^1a^*$ is essentially surjective. This is exactly the step where the mechanism fails.

Let's now identify exactly why $L^1a^*$ might fail to be essentially surjective. Let $U=W\setminus E$ and consider the span $\mathbb{P}^2\xhookleftarrow{x}U\xhookrightarrow{y}W$ that identifies $U\simeq \mathbb{P}^2\setminus S$ and satisfies $x=ay$. We find $x^*=y^*a^*$. Now $y^*$ preserves $1$-birational equivalences. Therefore, $$L^1x^*=L^1(y^*a^*)\simeq (L^1y^*)\circ (L^1a^*).$$ Since $x$ is $1$-dense, by \Cref{n dense equiv} we know that $L^1x^*$ is an equivalence. To show that $L^1a^*$ is essentially surjective, we want $L^1y^*$ to be conservative, and hence an equivalence by \Cref{fully faithful j_*}. However, since $y$ is only $0$-dense, $L^1y^*$ need not be conservative. In fact, $L^0a^*:\mathcal{H}^0(\mathbb{P}^2)\to \mathcal{H}^0(U)$ is essentially surjective, not $L^1a^*$. 

\subsection{$n$-birational descent of thom spheres}\label{4.5}
In this final subsection, we aim to show that there is an interesting way to produce lower-height birational motivic spaces from higher ones, namely by taking $\mathbb{P}^1$-loop spaces. To this end, we shall establish a Thom suspension and looping adjunction among the birational motivic homotopy categories of various heights:
\begin{prop}
   Let $E$ be a vector bundle of rank $n$ over a scheme $X$ smooth over the base. The suspension map $$Th_X(E)\wedge -: \mathcal{P}_{nis}(S)_\bullet\to \mathcal{P}_{nis}(S)_\bullet$$ sends $d$-birational local equivalences to $n+d$-birational local equivalences, and thus the delooping $\Omega _{Th_X(E)}$ sends $(d+n)$-birational local spaces to $d$-birational local spaces. In turn, this induces an adjunction $$L^{d+n}\Sigma _{Th_X(E)}: \mathcal{H}^{d}(S)_\bullet\leftrightarrows \mathcal{H}^{{d}}(S)_\bullet: \Omega _{Th_X(E)}.$$
\end{prop}
\begin{proof}
    Since the functor $$Th_X(E)\wedge -:=(\frac{E}{E\setminus0_X})\wedge -$$ is a left adjoint, it suffices to show that it maps $B_d(S)\cup \mathbb{A}(S)$ to $d+n$-birational motivic equivalences. For the case of $\mathbb{A}^1$-projections, consider the morphism $\mathbb{A}^1_Y\to Y$ in $Sm_S$. Then the map $$Th_X(E)\wedge ({\mathbb{A}^1_{{Y}}}_+)=\frac{E\times \mathbb{A}^1_Y}{(E\setminus0_X)\times \mathbb{A}^1_Y}\longrightarrow\frac{E\times Y}{(E\setminus0_X)\times Y}= Th_X(E)\wedge ({{Y_+}})$$
induced by the pointed projection ${\mathbb{A}^1_Y}_+\to Y_+$ is the Nisnevich local cofiber of 
    $E\times \mathbb{A}^1_Y\to E\times Y$ by $(E\setminus0_X)\times \mathbb{A}^1_Y\to (E\setminus0_X)\times Y$. Since both of the last two morphisms are motivic equivalences, the cofiber is so as well.
    
    Now, let $V\to U $ be a dense open immersion of height $d$. We have to show that $$\frac{E\times V}{(E\setminus0_X)\times V}\to \frac{E\times U}{(E\setminus0_X)\times U}$$ is a $(n+d)$-birational motivic equivalence. First, by Zariski excision, we have $$\frac{E\times V}{(E\setminus0_X)\times V}\simeq \frac{(E\times V ) \bigcup \big ((E\setminus0_X)\times U\big)  }{(E\setminus0_X)\times U}.$$ Now observe that $${(E\times V ) \bigcup \big ((E\setminus0_X)\times U\big)  }=(E\times U)\setminus\big (0_X\times (U\setminus V)\big).$$ So the map in question is the quotient by the
    equivalence $id:{(E\setminus0_X)\times U}\to {(E\setminus0_X)\times U}$ of the open immersion  $(E\times U)\setminus\big (0_X\times (U\setminus V)\big)\xhookrightarrow{} E\times U$.

    We claim that this last open immersion is $(n+d)$-dense. To see this, note that since $X$ is smooth over $S$, the closed immersion $$0_X\times (U\setminus V)=X\times (U\setminus V)\to X\times (U)=0_X\times (U)$$ is of codimension at least $d$, while the other closed immersion $0_X\times U\to E\times U$ is of codimension at least $n$, since $E$ is a vector bundle of rank $n$. And we are done. 
\end{proof}
  \begin{cor}\label[cor]{coneavue adjunction}
      Let $S$ be a Qcqs scheme. Then $\Sigma^n_{\mathbb{P}^1}$ takes 
$d$-birational motivic equivalences to $(n+d)$-birational motivic equivalences, and we have an adjunction, $$L^{n+d}\Sigma^{n}_{\mathbb{P}^1}: \mathcal{H}^{{d}}_\bullet(S)\rightleftarrows \mathcal{H}^{{n+d}}_\bullet(S) :\Omega^n _{\mathbb{p}^1}.$$ In particular, $\Omega_{\mathbb{P}^1}$ takes $(n+1)$-birational motivic local objects to $n$-birational motivic local objects.
  \end{cor}
      \begin{proof}
Since there is a motivic equivalence $\Sigma ^n_{\mathbb{P}^1}\simeq Th_S(\mathbb{A}^n_S)$, the result follows from the proposition above. 
\end{proof}

Let us now explain its significance. To do so, we first observe that the slice conjecture of Voevodsky [\cite{voevodsky2002possible}] can be restated as a combination of two statements: 



\begin{enumerate}
    \item $\Omega_{\mathbb{P}^1}(\mathcal{SH}/f_{n+1})\subset\mathcal{SH}/f_n $
    \item The induced morphism $f_{0/n}\Omega_{\mathbb{P}^1}\to \Omega_{\mathbb{P}^1}{f_{0/n+1}}$ is an $f_n$-equivalence.
\end{enumerate}

The corollary above is therefore the unstable analog of the first part of the (stable) slice conjecture. However, proving the unstable analog of the second condition is a nontrivial task.

Nevertheless, the above corollary yields an interesting tower of birational motivic spaces of varying heights, in contrast to the ordinary tower obtained in \Cref{bst}. Specifically, the following tower in $Pr^R$:
$$\cdots \mathcal{H}_\bullet^{3}(S)\xrightarrow{\Omega_{\mathbb{P}^1}}\mathcal{H}_\bullet^{2}(S)\xrightarrow{\Omega_{\mathbb{P}^1}}\mathcal{H}_\bullet^{1}(S)\xrightarrow{\Omega_{\mathbb{P}^1}}\mathcal{H}_\bullet^{0}\xrightarrow{\Omega_{\mathbb{P}^1}} \mathcal{H}_\bullet^{{-1}}(S)\simeq \mathcal{I}$$
And equivalently, the following tower in $Pr^L$
$$\cdots \xleftarrow{L^4\Sigma_{\mathbb{P}^1}}\mathcal{H}_\bullet^{3}(S)\xleftarrow{L^3\Sigma_{\mathbb{P}^1}}\mathcal{H}_\bullet^{2}(S)\xleftarrow{L^2\Sigma_{\mathbb{P}^1}}\mathcal{H}_\bullet^{1}(S)\xleftarrow{L^1\Sigma_{\mathbb{P}^1}}\mathcal{H}_\bullet^{0}\xleftarrow{L^0\Sigma_{\mathbb{P}^1}} \mathcal{H}_\bullet^{{-1}}(S)\simeq \mathcal{I}$$
\begin{rem}
Note that these towers run in the opposite direction to those in \S\Cref{bst}. As in \Cref{2.2.3}, it is an interesting question to understand the inverse limit of the first tower above in $Cat_\infty$. It seems that the objects of the limiting category will have certain interesting transfers. I do not yet have a characterization of this limit and hope to return to it.
\end{rem} 
\begin{rem}\label[rem]{Spq slice conj}
    Stating a similar version of the conjecture for the unstable spherical nullification of [\cite{asok2023p}, \S3] is rather easy. Indeed, by definition, for every $p\geq 0$ there is a tower in $Pr^L$ (of the $S^{p,q}$-nullified categories, defined as $\mathcal{H}^{p,q}_\bullet(S):=L_{S^{p,q}}\mathcal{H}^{\mathbb{A}^1}{(S)}$):
 $$\cdots \mathcal{H}_\bullet^{p+6,3}(S)\xrightarrow{\Omega_{\mathbb{P}^1}}\mathcal{H}_\bullet^{p+4,2}(S)\xrightarrow{\Omega_{\mathbb{P}^1}}\mathcal{H}_\bullet^{p+2,1}(S)\xrightarrow{\Omega_{\mathbb{P}^1}}\mathcal{H}_\bullet^{p,0}(S)$$   
 And the slice conjecture for this tower states that the canonical natural transformation $$L^{p+q,q}\Omega_{\mathbb{P}^1}\to \Omega_{\mathbb{P}^1} L^{p+q+2,q+1},$$ is a $L^{p+q,q}$-equivalence. By choosing a connected component of the base point, [\cite{asok2023p}, Corollary 3.1.21] can be used to restate this conjecture when the base is the spectrum of a perfect field. Namely, that for every $i\geq 1$  $(\pi_i\esc{X})_{-1}\to (L^{p+q+2,q+1}\esc{X})_{-1}$ is the universal strongly $\mathbb{A}^1$-invariant sheaf with a trivial $i$-fold $\mathbb{G}_m$-contraction.

 Also note that a simplicial version of the $L^{p,q}$ slice conjecture is known to be true for simply connected pointed motivic spaces (see [\cite{asok2023p}, Lemma 3.1.29 (2)]. Therefore, it is meaningful to state the unstable slice conjectures only for simply connected spaces.
\end{rem}
\begin{conj}[The unstable birational slice conjecture]\label[conj]{conj}
For a Qcqs scheme $S$, the canonical map $L^n\Omega_{\mathbb{P}^1}\to \Omega_{\mathbb{P}^1}L^{n+1}$ induced by \Cref{coneavue adjunction} is an $n$-birational equivalence on simply connected motivic spaces.
\end{conj}
    

\noindent\rule{\textwidth}{0.4pt}
 
\noindent{\large\textbf{\textit{Acknowledgments}:}} 

I thank my supervisor, Dr. Chetan Tukaram Balwe, for reviewing the initial draft. 